\section{Critical cutting, Palm comparison, and peeling}\label{sec:cutting}

This section connects the critical disks of Section~\ref{sec:proxy} to the
independently constructed sphere of Section~\ref{sec:continuum}, and then to the
exact map law of Section~\ref{sec:discrete}.  We first identify the critical
exploration together with every boundary contact.  We next identify the
counting intensity of sphere cuts, including its normalization.  Finally we
couple the exact peeling operations to these cuttings and obtain spatial
comparisons for all flag incidences, before and after exact volume conditioning.

Throughout the continuum approximation,
\(\kappa=\gamma^2\uparrow4\), \(r=4/\kappa\downarrow1\), and
\(a=Q_\gamma-\gamma\).  The disk area and boundary-length units are those of
Section~\ref{sec:proxy}.  The processes below describe the two actual sides of
a branch; their geometric identification is part of the proof.

\subsection{The critical branch and its contact relation}\label{cut:branch}

\begin{lemma}[The corner lifetime]\label{cut:lifetime}
Let the raw process have two independent symmetric Cauchy coordinates,
each with Lévy measure \(du/(2u^2)\), and kill it upon leaving the positive
quadrant.  Its Doob transform by \((x+y)^{-2}\), started from any
\((x,y)\in(0,\infty)^2\), has total mass one up to a finite lifetime \(T\).
For \(S_t=Z^L_t+Z^R_t\), one has \(S_t\to0\) as \(t\uparrow T\) and
\(\E T\le C(x+y)\).
\end{lemma}
\begin{proof}\leavevmode

The candidate length law is the Doob transform of two independent symmetric Cauchy processes (each density \(du/(2u^2)\)), killed outside the positive quadrant, by \(S^{-2}\), run to the corner. Here are lifetime details. At \(S=1\), twice the killed generator of \(S^{-2}\) is
\[
 \sum_{\ell=x,y}\left\{{\rm PV}\int_{-\ell}^\infty
 \frac{(1+u)^{-2}-1}{u^2}\,du-\ell^{-1}\right\}=0
\]
(each term is \(2\log(\ell/(1-\ell))+1/(1-\ell)-1/\ell\)).
Replacing \(-2\) by \(-1\) gives \(\sum(\log(\ell/(1-\ell))-1/\ell)\le-4\).
Scale homogeneity therefore gives transformed drift of \(S\) bounded above by a negative constant.
These calculations apply up to quadrant compact exits by Lévy optional stopping, including overshoots with the actual function value (zero outside; exits from pre-sides \(\ge\epsilon\) leave at least one side \(\ge\epsilon\), so the harmonic density stays bounded).
Thus there is a consistent probability law until leaving all compacta, with mean lifetime bounded by \(C S_0\) and \(S\) a supermartingale there.
Exits cannot run to infinity, nor leave through an individual axis while \(S\) stays positive: localized with \(S\) away from zero the density at inner side exits with \(S\) still bounded below is bounded, and on finite time windows a raw one-dimensional Cauchy process almost surely never hits a specified point by value or left limit (polarity, including jump pre/post values by the Lévy measure).
Jump exits outside are given zero weight.
Finally from arbitrarily small \(S\) returns to fixed large size before lifetime have probability bounded by the supermartingale estimate.
Thus almost surely \(S\to0\) at finite lifetime \(T\), and the localized laws exhaust unit mass.

\end{proof}

\begin{definition}[Side contours and the contact relation]
\label{cut:contact-relation}
Let $Z$ have the length law of Lemma~\ref{cut:lifetime}.  For each side
$\alpha\in\{L,R\}$, form an ordered contour by first following the original
boundary from target to source, with height increasing from $0$ to
$Z^\alpha_0$, and then following the side's length path from time $0$ to
$T$.  Fill every jump by its vertical height interval and give these
intervals positive summable traversal durations.  The resulting height
function $e_\alpha$ is continuous, is positive away from its two target
ends, and vanishes at those ends.

Two occurrences $u\leq v$ on this contour are directly linked if
\[
 e_\alpha(u)=e_\alpha(v)
       =\min_{w\in[u,v]}e_\alpha(w).
\]
The vertical interval at a jump is the aperture of that jump.  Its two
endpoints represent the same pinch; its open interval parametrizes the
jump-circle boundary away from the pinch.  At an ordinary real time there
is one tip copy on each side; at a jump the jumping side has a pre- and a
post-copy.  All tip copies of the same real time represent that time on
the trunk.  The source and target carry their usual endpoint copies.

The \emph{contact relation} is generated by the direct side linkages and
these tip and endpoint identifications.  It is a relation on occurrences,
so it does not presume that the corresponding embedded points coincide.
\end{definition}

\begin{proposition}[Critical cutting and the exact contact rule]\label{cut:critical-cutting}
Consider a critical field disk with a boundary-typical source and a target
at any prescribed strictly interior relative boundary length.  Its independent
symmetric \(\mathrm{bSLE}^{0,0}_4\) exploration admits a trunk-time
parametrization in which the left and right boundary lengths have the law of
Lemma~\ref{cut:lifetime}.  Conditional on the full length path, the jump
interiors are independent critical disks with the corresponding perimeters
and pinch marks.  At the compact-surviving deterministic times and selected
jump stops used below, the remaining disk has the marked disk law conditional
on the observed past length data.

The trunk is a simple continuous arc, all jump domains are Jordan domains,
and two contour occurrences have the same physical image exactly when
they satisfy the contact relation of Definition~\ref{cut:contact-relation}.  In particular no side class links two distinct real trunk times,
no class identifies distinct points of one outer or jump circle, and no
inserted positive circle visits the original boundary.  The full exploration,
including its ordered loop tours, is obtained as the limit of the subcritical
cuttings on compact surviving pieces, with the terminal remainder shrinking
spatially to the target.
\end{proposition}
The proof has three components.  We first construct subsequential limits
of the actual maps and traces.  We then prove that their contact relation
is exact.  Finally the Loewner driver identifies the limiting exploration
and its conditional disk laws.

\begin{lemma}[Compact-survival limits of the marked cuttings]
\label{cut:compact-geometric-limit}
Let good subcritical starting pairs converge to a positive critical pair,
and use the standard parent charts of Section~\ref{sec:proxy}.  Every
sequence has a subsequence on which the localized length-path couplings
extend jointly to the actual jump and post maps and the trunk.  The maps
converge uniformly on their closures and do not collapse; the trunk limit
is continuous and has no constant interval.  Small jump domains shrink
uniformly on compact surviving pieces, and the terminal remainder shrinks
to the target.  The positive jump boundaries are distinct simple
$\CLE_4$ loops.  The marked disk laws hold on every retained finite
past-length test, with the conditional scopes of
Proposition~\ref{proxy:marked-factorization}.
\end{lemma}
\begin{proof}


\proofstage{Localized convergence and stopping selections}

As \(r\downarrow1\), the subcritical strictly stable pairs converge locally in \(J_1\) before change of measure (the positive imbalance is \(O((r-1)^2)\), so even the compensation from large jumps converges without residual drift).
Hence the Doob laws on any pre-terminal compact path cylinder with null threshold boundaries converge with their weights, and the same is true jointly for the abstract jump and post disk laws on these pieces, by the preceding disk limits.
Here the (good subcritical) initial splits can tend to any desired critical pair in the open quadrant; positivity tests use strict interior compact ranges, or one conditions/exhausts by these.
There is no loss of probability by considering just such pieces up to post-jump times where \(S\) is small, because the candidate above has total mass one, reaches such times in countable threshold tests with probability tending to one (jumps are dense), and the \emph{entire remainder} at the corresponding subcritical jumps has diameter tending to zero in probability as the threshold for \(S\) tends to zero by \eqref{eq:original-11}.
This last assertion uses post disk law and a selectable stopping test bounded in time while the pre-stop path stays in a fixed positive compact set for each fixed cutoff test; one may also exhaust the selection by a bounded list and use disintegration at each first suitable jump.
Estimate \eqref{eq:original-11} conditional on the split/path supplies the upper costs with probability tending to one at rate a power of the small perimeter, independent of the number of steps for a stopping choice (the parent lower proxy test is made separately).
The remainder contains the terminal target in its closure.

For clarity these localized convergence arguments can be implemented by coupling along further subsequences of any input sequence: keep countably many finite time windows, side-range cylinder cutoffs and enumerated jumps bounded away from zero therein (choose thresholds off atoms), with marginal convergence on these compact-survival events.
One can place states at rational time horizons in a stopped path space with a default symbol if the path has already left the corresponding cylinder.
Cylinder masses and joint laws for all coarser initial pieces inside such cylinders converge by the bounded-density path transfer, including the stopped first-exit prefixes on smaller cylinders when needed.
Equivalently, for each accuracy first couple on finitely many disjoint compact-survival tests whose limit probabilities already sum to \(1-o(1)\) and which together continue through a selected tiny-\(S\) jump; then diagonalize accuracy.
Such disjoint tests are formed by rational bounds on duration and side extrema (including left limits), and testing the tiny-\(S\) entry using paths up to that bound (or first passing a jump threshold); priority for testing earlier entries can be imposed using only prefixes.
Failure to extend farther in the approximants never removes mass from a common prefix in a strictly positive cylinder.
This is a liminf/saturation use of compact convergence; it does not posit convergence of terminal durations at the singular corner.
Jointly bounded-coordinate field objects, paths and embeddings on the pieces need tightness, as follows.

In these localizations bounded-weight transfer works up to a first selected jump before a finite bound while pre-path and post-value stay in positive compact side ranges.
One way to see it just using deterministic horizon laws is to round that first jump time from above by dyadic times, requiring survival until then in a slightly expanded cylinder.
Right-continuity gives exhaustion at the actual stop for every counted compact-surviving path, and on the raw side the terminal density converges boundedly to the value using the jump-post position.
Equivalently stop the localized Doob law there.
For convergence across parameters, threshold/entrance and range tests up to the stop (including pre/post side sizes, the absence of earlier qualifying jumps etc.) can be chosen off all discontinuity boundaries in the raw limiting path; a fixed jump threshold has only finitely many candidates there before the time bound.
Thus this convergence of restricted finite measures applies also to their paths, and the disintegrated marked laws at those stops use the disk estimates uniformly in the selected lengths.
In particular one can use a \emph{single stopping selection} for each desired tiny bound \(\delta\) on \(S\), e.g. first jump above a lower threshold in absolute magnitude with small enough positive post-\(S\), before a sufficiently large fixed time and under pre- and post-compact survival bounds.
Choose these subsidiary thresholds (with strict margins, then allowing the compacts to increase for each fixed desired accuracy) so that existence has probability tending to one under the limiting transformed law, by jump density and \(S\to0\) only at terminal.
Conditional upper-chain bounds at this selected stop need not pay once per candidate.
The selected times in the candidate law eventually cover each initial compact piece as the tiny bound vanishes.
For post-map tests at other individual jumps on a fixed compact window, extra future success conditions need not be imposed when bounding bad-map probabilities: sum bounds at the jump stops themselves restricted to positive compact prefixes, controlling candidate counts there by the bounded-density and jump intensity estimates.
On projecting to a fixed matched post or cuts and length data up to their horizon, the marginal marked rule in any final subsequential limit follows by testing on these fixed compact-survival events in the approximants, whether or not maps from several different times were jointly kept.
Omitting outcomes failing the final tiny-\(S\) tests costs vanishing probability.
This use of exhaustion/couplings needs only convergence on initial segments; in expanding by excursions and passing to full curves the terminal remainder is controlled spatially, not by bounding its duration.

\proofstage{Tight closure maps and noncollapse}

Embed the original disk in \(\Delta\) in standard perimeter-normalized twice-typical coordinates (ignore the prescribed target for this choice if needed).
All individual jump offspring and post states on the tests above, mapped into \(\bar\Delta\) in their own corresponding standard coordinates, are tight in the uniform topology by \eqref{eq:original-11}; they are noncollapsed in subsequential limits.
Indeed their areas on inner compacta are tight with nonzero limiting mass, these are actual parts of the original area by interior covariance, and the parent measures converge to diffuse measures.
Interior field pullbacks therefore have the genuine charge-2 relation in the limit by univalent compact convergence, and lengths in each chart are the stated critical disk lengths.
Different jumped-away disk interiors remain disjoint and are avoided there by all past/future remaining traces (interiors persist by uniform compact convergence).
Taken collectively, small jump-disk maps on compact path intervals shrink in maximum diameter: use \eqref{eq:original-11} for the whole boundary with expected power sums of small jump sizes on the positive compact path tests bounded by stable-process sums of powers \(p>r\), tending to zero for a vanishing jump threshold.
The same argument can transfer high-probability short-arc tests on countable jump-post lists; moment orders in \eqref{eq:original-11} can be chosen arbitrarily high.
Shrinking here and below takes place jointly using Lemma~\ref{proxy:lower-cost} once in the ambient chart, then short-cost upper tests.

\proofstage{Oscillation of the actual trunk}

We verify tightness of the actual \emph{trunk in trunk time}, not just of the domains.
On compact-survival intervals partition at successive jumps of size \(>\delta^R\), with fixed \(R>1\); include the initial state and, if stopping at a selected jump, end the last interval there.
For each fixed compact cutoff under consideration, with probability tending to one as \(\delta\to0\) there are \(O_{\mathbb P}(\delta^{-Rr})\) such intervals and both side lengths oscillate by \(<\delta\) between the initial post-state \(a'\) of each and just before its last jump \(b'\) (different constants in front of \(\delta\) would work too).
For example under the base laws the longest gap before a bounded horizon is at most \(O(\delta^{Rr}\log^2(1/\delta))\) with high probability.
Compensated smaller increments have variance rate \(O(\delta^{R(2-r)})\) and jump bound \(\delta^R\), hence exponential maximal bounds by Lévy exponential tilting on blocks of that length; the remaining skew drift has size at most \(O(\delta^{R(1-r)})\) uniformly.
Transfer on the cylinder.
By \eqref{eq:original-11}, all length-\(O(\delta)\) arcs on the post and terminal jump frontiers under test have simultaneously vanishing image diameter in probability.
For jump disks of sizes \(s\) possibly much smaller than the side lower bounds, use the whole-boundary estimate if \(s\lesssim\delta\), otherwise \eqref{eq:original-11} with relative size \(C\delta/s\), bounded perimeter factor and the grid union.
The polynomial number of tests costs nothing with sufficiently high moments.

In the canonical starting disk at \(a'\), near-full prefixes (leave \(C\delta\) gap to the tip) on both sides are then untouched until \(b'-\). The non-inherited arcs of the boundary of \(D_{b'}\) are short, with the following jump qualifications.
\begin{itemize}
\item At a positive terminal jump one inserted aperture may be large; join back \(C_{b'}\) along that aperture, whose open seam outside the pinch has a collar on both sides and was not encountered earlier.
\item At a large negative jump the hit is on an untouched prefix; the long segment of \(\partial C_{b'}\) from that hit back toward the old tip (ending on the prefix) is inherited, and all the rest of that detached boundary has length \(O(\delta)\). At a small negative jump the whole detached boundary is small.
\end{itemize}
In all cases the omitted \(O(\delta)\) arcs at \(a'\), all the changed short arcs at \(b'\), and both tips are inside a parent ball of small radius tending to zero under our tests: they link by common inherited prefix endpoints and common pinches.
The descriptions follow from the strict stack rules of Lemma~\ref{proxy:stack-rule}, choosing margins off the fluctuating side lengths before the last jump; in the borderline case of a negative loss within that margin use only shorter prefixes and the small jump bound.

No intermediate macroscopic finger of the trunk or removed material (except the cell at the last jump itself) can then lie outside this ball, enlarged a little.
Indeed in the open canonical disk of \(D_{a'}\), consider each component outside its pullback under the parent embedding.
The boundaries of the post and terminal jump domains there have no interior portions except, in the positive case, the shared open seam just described.
Thus a component is wholly retained in their union (including this seam), or would be wholly discarded.
Any discarded component must have an access adjacent to an untouched/inherited old boundary point mapping outside the ball: otherwise its whole boundary in the closed disk maps inside the closed ball, impossible by the maximum principle.
At such a boundary point the post or terminal jump domains have genuine inherited collar access on this starting-disk side, contradiction.
Boundary contacts outside the ball at earlier times are likewise excluded by prefix protection.
The interiors of post and last-jump domains are never touched on the trunk interval and the large positive seam is first reached at the last tip only.
This proves the claimed tip oscillation up to and including the terminal pinpoint on each interval.
For fixed mesh parameter successive large jumps are separated in time in the localized limit.
Thus these oscillation controls prove tightness with continuous limits under the \(J_1\) path couplings of lengths.
At the end, after tiny-\(S\) localized stops as above, all remaining tips and jumps are confined near target.

Pass now to joint limits on these pieces, using time homeomorphisms as needed for \(J_1\).
We have a continuous limiting trunk \(P\) up to and including \(T\), continuous closure maps with inherited disk laws at jumps, and a full continuous path obtained by interspersing the positive loops, traversed in boundary-length order of the corresponding orientation.
This last path also has uniform convergence modulo time change: give positive auxiliary durations to the positive jumps summably in their countable matched order, or match just finitely many large ones at each accuracy, using small diameters and trunk continuity for the rest.
The trunk avoids jump-disk interiors by approximation.
Every positive loop here is a distinct simple outermost loop of a genuine \({\rm CLE}_4\) in \(\Delta\).
Indeed consider also the parent non-nested CLEs through their domains at countably many interior tests.
These domains converge in kernel law to those of \({\rm CLE}_4\): by the Sheffield--Werner Brownian loop-soup construction \citep{SWcle} they may separately for this assertion be coupled increasing at each interior point as soup intensity increases, and the Schramm--Sheffield--Wilson conformal-radius law \citep{SSWradii} is continuous up to \(\kappa=4\).
Thus the increasing kernels exhaust the final interiors (strict radius monotonicity); prescribed interior points lie in loop interiors almost surely.
This is tight joint data before any localization.
Each positive jump domain contains a rational interior point stably, so equals in its interior limit the corresponding kernel, by its uniformly convergent univalent map, and distinct ones have disjoint interiors.
Its boundary map thus parametrizes the simple CLE boundary homeomorphically.
Consequently \(P\) has no constant intervals, since distinct disjoint simple loops attach on any such interval by positive-jump density.


\proofstage{Consistency of the limiting data}

The compact extractions can retain all the path data consistently.
Localized extractions as above need only tightness of joint embedded data; for the post/cut \emph{laws} they use the marked rule tested up to each piece's own horizon or stopping jump (not a claim that post surfaces at different times are conditionally independent).
For finite unions of compact path tests with varying horizons, intersections are again tests on a longest compact-surviving path and inclusion-exclusion gives path probabilities.
Thus one can keep the path variables consistently in these joint extractions.

\end{proof}

\begin{lemma}[Exact contacts in a subsequential limit]
\label{cut:exact-limit-contacts}
For each limit in Lemma~\ref{cut:compact-geometric-limit}, the contact
relation of Definition~\ref{cut:contact-relation} is exactly equality of
physical images.  No class links two distinct real trunk times, and at a
real time at most one nontrivial side linkage occurs.  The trunk is simple,
all jump domains are Jordan domains, and no positive jump circle visits
the original boundary.
\end{lemma}
\begin{proof}
\proofstage{The auxiliary prime-side disk}

The following quotient argument identifies \emph{all} contacts of these maps; it is important not to infer an inverse welding rule just from interior kernel convergence.
Make an auxiliary clean closed disk with a simple cross-diameter parametrized by \([0,T]\) from source to target.
Put at every jump a small closed round disk tangent at its time on the diameter, on the jumping side, otherwise away from original boundary and diameter, with diameters tending to zero and disjoint from all the others (choose them successively).
In the remaining open strip region of each side use its closed prime disk.
Its prime boundary goes from the target along the old boundary back to source (label by height increasing from 0 to \(Z^i_0\)) and then along the diameter to the target detouring around each tangent circle on that side, labelled with height \(Z^i\) using linear interpolation from pre- to post-height over the detour.
This is a continuous excursion once the detours are given circle parameters.
The prime-to-physical map on this side folds only the two copies at each pinch on the diameter.

For detail, the region is simply connected (the excluded tangent disks join the exterior) with locally connected boundary by their vanishing sizes.
In its boundary tour all open outer arcs and circle arcs away from pinches have a unique access. A diameter point which is not a tangency has a basis of separating crosscuts tending to it enclosing all near accesses: for any small neighborhood first isolate it between diameter points away from the finitely many large disks, connect across slightly above the diameter avoiding those disks, and approach at the two chosen points (not tangencies) with short paths staying on this side by detouring around smaller circles.
Equivalently paths in the unobstructed half-neighborhood may be adjusted around intervening small tangent disks of diameter tending to zero; one can use their outer semicircular boundaries with arbitrarily small clearance, as each positive-height compact piece stays away from all but finitely many, and approach the line by successive such depths.
At a tangency the same argument on the two separated sides of its excluded disk gives one access each.

More explicitly for a diameter point off all tangencies, at small horizontal scale one can truncate at a fixed nearby horizontal height and delete finitely many disks reaching that height.
In smaller rectangles centered at the point all interfering disks are as small as desired.
To build barriers, choose two tangency-free base points on either side with upward access obtained by connecting a sequence of free interior points converging to each one, sliding horizontal/vertical segments into the exteriors of interfering disks (all sliding distances vanish as one approaches the point), then simplifying to a crosscut arc.
Since such disks are disjoint, each detour at positive height along one stays locally a positive distance from all the others.
This also works separately on each flank of a fixed tangent disk at its pinch.
These crosscuts may tend to the point with the bounded side containing all sufficiently near accesses in the stated sector.
This gives the ordered prime circle with the indicated folds, also by cyclic order of the crosscuts and the open detours.
In the prime-side geometry description, another precise way to check uniqueness of near access used there is the following.
At a diameter point not a tangency, any two interior points tending to it can be joined by paths of vanishing diameter: start with interior segments near it at positive height, replacing portions hitting the finitely many tangent disks reaching those portions by detours along their outer arcs avoiding the line.
These arcs at positive height have clearance from all other disks and can be slightly pushed off; sizes of interfering disks go to zero.
At a tangency the same works separately in each flank (make the initial paths in the flank outside the one fixed disk before adjusting for all the others).
Two genuinely inequivalent prime labels over the point can be separated in the round prime disk by a crosscut with endpoints mapping away from the point (no boundary interval is constant), thus cannot have this local joinability.
This verifies the access counts including in the case of densely spaced jumps.

We compactify outside the auxiliary closed disk to a sphere.

On each prime side circle impose the direct linkage relation of Definition~\ref{cut:contact-relation}. Collapse the closed convex hull of each such class in the round prime disk. Here are properties holding almost surely for our Cauchy paths:
\begin{itemize}
\item No class on a side contains two distinct real diameter times (count pre/post at a jump as the same time). Apart from the trivial target/source conventions, at any real time at most one nontrivial such linkage to other labels occurs among the different side/time copies (a single nontrivial side class, which may itself have several labels).
\item Thus in the two-sided physical auxiliary disk the collapsed sets, combining across folds and the common diameter, are nonseparating compact connected sets making an upper semicontinuous decomposition, with no identifications on any single outer or tangent boundary circle between distinct physical points of that circle.
Equalities between different contour occurrences away from the fold/identical-time conventions are just the indicated direct stack linkages to the circles, arc labels or trunk time (when at a folded/time tip, use the incident copy making the link).
\end{itemize}

\proofstage{Record-range polarity}\phantomsection\label{cut:range-polarity}

For the first assertion use Cauchy fluctuation theory on intervals strictly before terminal via localized absolute continuity.
At a rational split between two linked diameter times their common level would be in both closed ranges of the descending record heights from that split in the two directions, relative to the split level (vertical jump \emph{endpoints}, not the filled-in jumps, for real diameter times).
These independent ranges for raw Cauchy paths are polarized apart off their origin: the descending ladder heights, like the descending ladder times, have the range of a stable subordinator of index \(1/2\).
We use the standard Cauchy ladder-process identification from Lévy fluctuation theory \citep[Chapters~VI and~VIII]{BertoinLevy}.  For the range estimate, Taylor--Wendel's exact Hausdorff-gauge theorem and its stable-subordinator interpretation imply zero \(\mathcal H^{1/2}\) measure for each bounded closed range piece; adding the countably many left limits does not change this conclusion \citep[Theorem~1 and Section~3]{TaylorWendel}.
Such a range avoids any independent translated or reflected range away from a deterministic starting value: hitting an interval of width \(\epsilon\) on a compact positive-axis subrange has probability \(O(\sqrt\epsilon)\) by the stable overshoot law \citep[Lemma~1]{TaylorWendel}, hence sets of zero \({\cal H}^{1/2}\) there are polar.
At the split itself the path immediately goes below, so the common split height does not qualify.
This proves no distinct-diameter chord.
Similarly there is no time simultaneously sustaining a real-time record interval (in either direction) for the two sides, using the descending ladder \emph{time} ranges viewed toward such a time from rational points independently for each side.
At a jump, on the same side the pre-value cannot sustain a linkage backward, nor the post-value forward (immediate oscillations by jump Palm/independent increments); only one of the reverse choices can extend over the jump depending on sign.
The opposite side has no record there.
At source the path immediately dips on each side, and target has zero height, excluded everywhere else.
This proves the assertion.

\proofstage{The closed decomposition and its quotient}

A side hull is compact and projects injectively in the auxiliary disk (the two folded endpoints of a jumping side have different heights); side hulls are disjoint before folds by noninterlacing and closed as a relation by continuity.
After folds any merging is only at a \emph{single} diameter point for a cluster of hulls (each hull has at most one such physical point); thus chain length is uniformly bounded, preserving closedness.
Each set is nonseparating on the sphere, being a planar embedded contractible compact convex set or finitely many such attached at one point (e.g. planar Alexander duality).
Injectivity on a tangent circle now follows by monotonicity of its aperture heights and absence of two different nontrivial hulls merged through another time.
An interior aperture label also cannot tie back to its own pinch via folding, by the same observation and the unequal heights.
Injectivity on the original outer circle follows likewise (no cross-side nontrivial linking at one time, including the endpoint conventions).
All points of the side prime disk already belong to the union of hulls: for any two nonequivalent boundary labels a level chord strictly separates them, taking endpoints of a connected component of the path above an intermediate height.
If a point in the round disk were not covered by the closed union, for each class hull keep the closed circular segment outside it on this point's side (containing the point).
The intersection has interior, with extreme points only on the circle (bounding chords do not cross; any active chord in the closed limit stays uncut in the disk, so its relative interior is not extreme).
All these extreme points are pairwise not strictly separable by any level chord, hence equivalent, contradiction.
These facts give the claimed decomposition; interior points of tangent disks and outside are left untouched.
By Moore's decomposition theorem \citep[Theorem~22, p.~425]{MooreDecomposition},
the quotient is a sphere.  To apply its plane formulation, take an untouched
exterior singleton as infinity.  Every other class is a bounded plane
continuum; its complement remains connected after removing that singleton.
The restricted decomposition is upper semicontinuous.  Moore gives a plane
quotient, and restoring the singleton in the compact Hausdorff quotient gives
its one-point compactification, a sphere.

\proofstage{Transfer of all stack linkages}

Every stack linkage just collapsed really has common image in the cutting limit above.
Strict ties from an open outer/positive aperture label into an open loss aperture have this property by persistence and the strict subcritical comparison, also with non-strict intervening equality by perturbing both heights a little downwards within the apertures.
For a tie to or from a real time use the same perturbation into a strict prefix at a large-jump partition time, within \(O(\delta)\) of its top height, observing that (by taking a sequence of meshes and the preceding simultaneous arc estimates) all such short arcs have vanishing image size in the limit on compact pieces.
More explicitly at a continuous time in that coordinate use the preceding post-state \(a'\) nearby, whose height differs by \(O(\delta)\); transmission from an earlier aperture extends to \(a'\), or forward transmission from \(a'\) to a later aperture holds, on lowering by a sufficiently large multiple of \(\delta\).
The interval from \(a'\) to the continuous time also has oscillation \(O(\delta)\).
Partitions can be assumed to have these properties and arbitrarily small time mesh eventually along further subsequences, pathwise. \(J_1\) comparison with the approximating path then transfers the strict labels, by first fixing mesh margins/endpoint perturbations then passing to that subsequence limit.
If at a jumping time, use its own post-state: either the linked height is the post-value and just below is near its top, or there is a forward tie from the pre-value over a positive insertion, in which case just below is adjacent to the pre-height end of the inserted arc, also mapping to the tip. A pre-jump backward record does not occur here. A nontrunk earlier occurrence tying to a later label/time has to be on an outer increasing arc or open positive jump; a later nontrunk occurrence on the detoured path in a linkage has to be on an open loss interval.
Thus the perturbations suffice for all cases by aperture endpoints mapping to the tips and positive compact exhaustion.
Source and target have only their automatic endpoints.
This proves the image claim simultaneously (the pathwise strict-margin transfers allow arbitrary labels on good samples).

\proofstage{Excluding additional contacts}

Map the quotient sphere to the actual compactified disk using \(P\), the inherited outer boundary, and the jump maps with their boundary labels (using any homeomorphic abstract identification of each inserted auxiliary disk with the prime disk of its jump matching these labels); extend over the exterior homeomorphically.
The map is continuous: it descends from the closed exterior, diameter and all tangent disks including their interiors taken together (a closed subset hitting all classes); continuity there follows from shrinkage at accumulating jumps.
It is a local homeomorphism with singleton fibers at the images of all jump interiors and the exterior, by noncollapse/disjointness and avoidance by trunk and other boundaries.
These regular points are dense on the \emph{source} side by jump density and the hull description.
Here is why no extra contacts form.
Degree on the sphere is nonzero by any of those singleton regular fibers, so the map is onto.
Moreover any disconnected fiber could be separated into two compact portions with open neighborhoods whose boundaries avoid the fiber.
Local degree on each neighborhood would be constant on a common ball about the value.
By regular point density and their singleton fibers, each degree must be nonzero, which now contradicts existence of nearby singleton fibers.
Thus the map is monotone.
But each remaining fiber is zero-dimensional: it lies on the countable union of outer circle, diameter, and tangent circles embedded injectively as just proved in the quotient; its closed intersection with each contains no interval, by absence of trunk plateau and by nonconstancy on intervals of continuous univalent boundary maps (reflection or boundary uniqueness).
Each intersection is a zero-dimensional closed subspace of the fiber.  The countable closed sum theorem for separable metric spaces \citep[Theorem~1.3.1]{EngelkingDimension} therefore makes the entire fiber zero-dimensional.
Hence every fiber is a singleton.
This proves homeomorphism and the exact no-extra-contact stack rule, in particular simple trunk, Jordan offspring including the cut-offs, and no original boundary visit by an inserted positive circle.


\end{proof}

\Needspace{5\baselineskip}
\begin{lemma}[Identification of the driver and conditional disks]
\label{cut:limit-identification}
The limit in Lemmas~\ref{cut:compact-geometric-limit} and
\ref{cut:exact-limit-contacts} is the independent symmetric
$\mathrm{bSLE}^{0,0}_4$ on the parent critical disk.  Its length path is
the law of Lemma~\ref{cut:lifetime}, in the clock recovered by positive
jump counts.  Conditional on the full length path, its jump disks are
independent disks with their stated perimeters and pinch marks.  At each
tested surviving deterministic time or selected jump stop, the remaining
disk has the marked law conditional on the past length data at that time.
\end{lemma}
\begin{proof}
\proofstage{Loewner hulls, excursions, and the intrinsic clock}\phantomsection\label{cut:driver-clock}

Now use the upper half-plane Loewner chart (source and target at 0 and infinity), normalized also by any nondegenerating parent-coordinate convention.
Away from the terminal point the capacity of full-trace initial segments converges under the uniform path couplings above.
Indeed the limiting trace avoids infinity up to every pre-terminal time by the stack rule and the interior location of the positive loops.
Bounded hulls then have kernel convergence as viewed from infinity by uniform trace convergence (paths in the unbounded component persist on compacta, and interior boundary points on the trace prevent a larger kernel), which gives hydrodynamic convergence and convergence of capacity.
Capacity is continuous and strictly increasing in the expanded parametrization (include positive-loop intervals, no inserted holds).
For any inserted positive loop at time \(s\), its post disk map has the whole loop as a boundary subarc, and an open true collar toward target \(y\).
Indeed actual untouched prefixes on both original arcs of fixed short lengths near \(y\) persist strictly through \(s\), with exact length labels; the limiting post map on that open arc is onto an original boundary interval, hence reflects there.
Its interior is unvisited so is accessible from \(y\).
No part of the positive loop was visited \emph{before} its attachment: a strict insertion label cannot tie backward over its jumping time, and the tip there cannot tie backward by any side.
During the loop traversal, forthcoming distinct interior-of-parent loop points are therefore still reachable (in the closure of the post domain and not in the earlier trace).
This gives strict hull growth on loop intervals, hence everywhere by positive-jump density.
These checks can also be used for continuity of the unbounded kernels at any pre-terminal path time.

Jointly keep the driver pair of the independent bSLE, tight and converging to the standard symmetric critical driver with its excursion intervals.
In detail \(X=(W-O)/\sqrt\kappa\), driver minus force-point location divided by \(\sqrt\kappa\), is the symmetric signed Bessel of dimension \(\delta=3-8/\kappa\), and \(O=-(2/\sqrt\kappa){\rm PV}\int dt/X\) in chordal Loewner time, with symmetric cutoff (Sheffield's side-swapping convention).
The nonzero excursions trace the discovered loops.
As \(\delta\uparrow1\), use common Brownian motion in scale coordinate \({\rm sgn}(X)|X|^{2-\delta}\), with elapsed bSLE time per Brownian time at height \(u\) of the motion equal to \((2-\delta)^{-2}|u|^{-2(1-\delta)/(2-\delta)}\) (speed-scale construction without holding).
Clocks converge compact-uniformly by Brownian local time.
Principal-value integrals likewise converge using the difference of local times at \(u,-u\), uniformly Hölder at zero of any fixed order less than \(1/2\) on compacts; the remaining kernel has power \(-(3-2\delta)/(2-\delta)\) there.
This proves the asserted marginal convergence jointly with zero sets/excursions, and independence of the limiting critical driver from the parent field given the endpoint prescription.
Here Loewner time uses \(\partial_t g=2/(g-W)\), i.e. half-capacity.

The limiting hull domains throughout are driven by that pair's \(W\), by driver and time convergence (standard chordal Loewner kernel convergence, or solve the inverse flow on compacta).
Total Loewner time is infinite, since otherwise the continuous driver on a bounded interval could not generate the arbitrarily far interior loop points as the trace tends to target. A continuous capacity-parametrized path giving these strictly growing hull domains is unique: if two such paths separated at a tip time, the domains at slightly later times are strictly smaller; their boundary in the preceding open domain is nonempty and must be on both short incremental traces, impossible in disjoint neighborhoods.
Thus by the bSLE continuity results the limit is exactly symmetric \({\rm bSLE}^{0,0}_4\) (i.e. the symmetric chordal \({\rm SLE}_4(-2)\) construction without added local-time drift).
Excursion intervals and their signs match, by strict capacity increase and speed-scale excursion convergence (nonvanishing loops have nonvanishing limiting time interval there, and vanishing loops away from infinity only vanishing growth).
In particular we can use the MSW critical GFF/CLE-percolation coupling: in a Dirichlet disk with constant boundary level and independent GFF \(J\) with that constant added, this bSLE has the realization tracing signed CLE\(_4\) loops on first encounter with the central level line (boundary-height level, law \({\rm SLE}_4(-1;-1)\)) as trunk.
The exploration/branching Markov coupling gives fresh CLEs in unexplored cut domains, and in this realization Dirichlet GFFs there at the unchanged constant in negative regions and at increments \(\pm2\lambda\) in positive loop interiors (\(\lambda\) the Schramm--Sheffield half height-gap), with future branches using the corresponding level lines.
This uses the Dirichlet disk coupling (pull back by conformal maps), not an assumption about geometry of LQG cuts.
See also the level-line/local-set construction in the MSW proof. \citep[Proposition~5.3]{MSWTrunk}

Jump disks really are exactly the indicated complementary domains, including the negative cuts in this unclocked description.
All negative boundary labels link to earlier labels or tips by time of the cut (trace the height from earlier on that side), so their Jordan boundary lies on the trace and/or original boundary by then.
At any matched jump time \(s\) the post-domain map's interior is in the accessible target component as seen above (the collar argument works for either jump), and contains interiors of all strictly later jump disks: by uniform post-map convergence it contains them at least in closure, and its boundary lies on the earlier trace/original boundary by approximation, hence outside their interiors.
Thus later disks were not removed prematurely.
Conversely there are no additional detached open components by the homeomorphism: the jump interiors are dense, and earlier ones already have their full boundary obstacle by their event times.
More explicitly the terminal complement in the original open disk of the trunk and positive-loop circles is precisely the union of all jump interiors by the quotient representation.
This identifies disconnections at negative jumps and their pinch tips, and the target post-domain interior equals the actual remaining component (also by the kernel theorem from infinity since the convergent map has interior there).
Orientations and inherited length labels are joint physical ones from the limit.
At a negative jump the side can also be read at the tip of the approaching trunk in the directed quotient disk.
The truncating time change on this bSLE is now intrinsic at least in the almost-sure sense needed here: lengths of its encountered positive loci from the \emph{inside}, as field disks, are measurable from the parent and the trace by the disk length identification.
Counting the large ones recovers trunk time by Cauchy absolute continuity, simultaneously along the branch.
Similarly signed changes in the limit here are read from their assigned jumps (symmetric compensated Cauchy summation), or measured on the post frontiers.
The jumping side of a detached interior is intrinsic: the directed limiting trunk is a simple diameter-type arc with the two parent sides specified (one may close it externally by an arc from target back to source to define its sides even with contacts to the original boundary).
By the quotient homeomorphism this agrees in actual disk coordinates with the side assignment.
Cut and continuing fields in achieved domain interiors are the actual restriction pullbacks, since on all interior compacta univalent map convergence transports the tested field distributions (localize images to compacta in the parent).
Boundary lengths needed here agree with their field-as-surface-function versions by the disk marginals.
Thus to realize the one-branch law with the prescribed decoration one uses the limiting \emph{pair} driver/excursion data, independent of the parent field conditional on endpoints, and the MSW joint trunk/encountered-loop coupling with exactly that symmetric bSLE excursion law.
Curve and trunk/discovery order are fixed by these driver/hull/excursion data as above.
One may condition that ordinary geometric coupling on them to add its unobserved decoration, without retaining any tentative coupling with unhit CLE kernel tests from the subcritical extraction.
The marginal of \(J\) used to generate this branch is the ordinary independent constant-boundary Dirichlet one conditional on the parent.
There is no residual freedom in the matched lengths/coordinates other than harmless chart conventions. \citep[Proposition~5.3]{MSWTrunk}

\proofstage{Remaining disks at deterministic times}

At any ordinary deterministic compact-surviving trunk time the same remaining-disk rule holds.
For instance take nearby later jump stopping times (first jumps over vanishing thresholds after that time), exhausting with positive compact bounds.
Their actual critical post maps just identified are tight uniformly up to closure even in this varying list: for each threshold the laws arise by the matched subcritical limits and \eqref{eq:original-11} and the single parent lower test with uniform constants, conditional split bounds as needed.
Normalized maps are noncollapsed by area; as times decrease, kernel convergence as viewed from target and convergence of traces identify the interior at the deterministic time, using an untouched target collar as before to exclude a trapped normalization.
Tips and targets converge to the appropriate prime ends: target remains untouched, and at deterministic time almost surely no record linkage extends to other labels, hence the point is new and there is only one access there in the target component.
Indeed an inequivalent access would again give a joining interior crosscut from and to this point enclosing a pocket disjoint from the prior trajectory from the outer start up to its terminal visit, forcing a constant boundary interval.
Interior fields therefore identify the abstract post law and its marks by the converging split disk laws; joint testing with completed marked cuts and lengths gives the claimed conditional independence statement by localization and right-continuity.
One can approximate general length stopping pieces analogously when required, or work with jump and dyadic times.
The lengths/clock and contact law just proved hold for an independent MSW branch in an actual critical disk, not only an unspecified geometric realization of the Cauchy data, by the above driver identification and recovery.
In particular one can resample its geometry by the \(J\) coupling instead; countably many Markov branches can be routed together, exploring onward in detached domains with relative targets prescribed from transmitted boundary labels or freshly length-sampled, under the same one-branch rules.
Conditioning for lengths and offspring here retains only the abstract previous length data/pinch prescriptions as tested in the factorization; embedded data need not factor.

Product cut-disk laws given the full parent length path follow already by compact survival exhaustion: for finitely many selected cuts observed on a surviving bounded prefix, test the deterministic-horizon formulas also against later bounded surviving path events.
Finite combinations (including later nonsurvival by subtraction) and monotone class then give the claimed cut marginal even conditional on the whole length path.
This does not claim such independence for \emph{post} disk fields given their future path.


\end{proof}

\begin{proof}[Proof of Proposition~\ref{cut:critical-cutting}]
Approximate the prescribed starting pair by good subcritical pairs as in
Lemma~\ref{proxy:jump-times}.  Lemma~\ref{cut:compact-geometric-limit}
produces the limiting actual maps, traces, and marked disk laws.
Lemma~\ref{cut:exact-limit-contacts} identifies every contact, and
Lemma~\ref{cut:limit-identification} identifies the exploration as the
specified independent critical branch.  The shrinking terminal remainder
extends the compact-survival comparison to the complete ordered loop tours.
These conclusions give all the assertions of the proposition.
\end{proof}

\begin{proposition}[One auxiliary field for compatible branches]\label{cut:branching-fields}
The critical branches can be realized as full signed-CLE explorations of
one constant-boundary Dirichlet GFF \(J\), independent of the Liouville input.
Branches with the same source and compatible transmitted targets agree until
the targets separate, and continue in the resulting components with the
signed exploration Markov law.  Negative complementary components have the
unchanged constant boundary value; positive loop interiors have the signed
increment \(\pm2\lambda\) and retain their independent Dirichlet fillings.
These statements concern the full exploration, including encountered loops.
They allow countably many compatible continuations and preserve the abstract
marked-disk factorization stated in Proposition~\ref{cut:critical-cutting}.
\end{proposition}
\begin{proof}\leavevmode

We detail the consistency here with multiple targets and \(J\).
We use the signed CPI version of MSW in a constant-boundary Dirichlet disk: jointly with the signed non-nested CLE the central level line plus tracing of its encountered loops is the symmetric bSLE exploration.
Unencountered non-nested signed loops in negative cut components (and in a continuing target component at a branching restart) have the fresh signed CLE Markov law; when branching to two still-transmitted boundary targets the ordinary bSLE exploration coupling follows a common past up to separating them, then restarts toward the targets in their respective components.
Here only full explorations including the loop excursions are considered for the cut components, and separation on the original boundary is at a negative disconnection with a common tip, not during a partial positive-loop excursion.
These uses are in ordinary disk conformal type, with countably many ordinary endpoint choices independent of the current Dirichlet GFF before exploration (one can condition initially on independent Liouville data).
The statements use the CLE-percolation/branching domain Markov rule, not a claim that a slit disk obtained by uncovering the \emph{trunk alone} has constant data. \citep[Proposition~5.3]{MSWTrunk}

Here is also a reconstruction explanation for using this coupling with Dirichlet fillings and level continuations.
In the constant-data CLE\(_4\)/GFF coupling the signed carpet is the two-valued thin local set with offsets \(\pm2\lambda\); conditional on it, fields in loop interiors are independently Dirichlet at their labels. A central level line trunk to an ordinary target is a deterministic local set of \(J\) staying outside all these interiors.
Thus it does not bias their fillings even on conditioning on it as well as on the signed loops; indeed it is determined by those signed loops.
To see the locality implication, for a fixed finite union of test domains \(O\) compactly in the starting disk, on avoiding its closure the level line is measurable from the field off its zero-Dirichlet projection in \(O\) (deterministic-local-set avoidance property; its simple trace fixes order).
Conditional on the signed loops, use increasing such test sets with closures inside their interiors and union exhausting all the interiors, chosen from countable libraries.
Conditional avoidance there has probability one.
In that conditional Gaussian Dirichlet-filling law the observed outside projection uses only the labels and the complement of the same interior projection in the filling GFFs: projections onto test functions supported inside \(O\) are just the corresponding Dirichlet filling projections since boundary shifts there are constant.
These decreasing complementary Gaussian sigma-fields have trivial intersection conditional on the signed loops, by exhaustion/density of the interior Dirichlet spaces in the sum of the loop Dirichlet spaces.
This proves the assertion (and the full bSLE path/traversal is then fixed by its trunk and the oriented loops).
Consequently in negative complementary domains with the fresh signed-loop law the field is the corresponding constant-data Dirichlet GFF, and positive loop interiors also retain their prescribed independent fillings.
Indeed restrictions to unencountered loop interiors agree with the independent fillings shifted by their signed labels in an ordinary signed CLE/GFF of each new domain, by the signed Markov rule and the preceding conditional independence.
No discrepancy on the remaining carpet inside the open domains is possible: its neighborhood volume decays as a positive power on compacta a.s. (exponential tail of the single CLE conformal-radius drop at interior points and Koebe, integration and dyadic Borel--Cantelli), while both the ordinary disk GFF and the restriction of the original GFF there have arbitrarily small negative Hölder order.
Cutoffs on fine squares meeting the carpet with smooth rescaled derivative bounds thus have vanishing pairings with the difference.
This works jointly across the cut domains.
Alternatively this is the constant-data complementary-region rule of the MSW full level-set exploration. \citep[Proposition~5.3]{MSWTrunk}

In particular continuations can be taken in the same \(J\).
At a target fork one can use the ordinary target-consistent signed bSLE coupling; its branch with either target and the parent signed loops has the one-target CPI joint law, so must agree with the corresponding deterministic level-line exploration in this field.
In a detached negative disk branching into it from the common tip uses the ordinary signed bSLE in that disk, hence by the same joint-law determinism agrees with the corresponding level exploration there.
Equivalently, after a full branch redirect to swallowed targets by these rules in its fresh components.
Only targets on open surviving/transmitted intervals, not ambiguous corner targets, need be redirected.
The signed geometry and recursive continuations can be conditioned also on the initially independent Liouville input throughout these applications.
For length laws one nevertheless conditions only on the tested abstract path/mark data to get fresh field disks, not on embedded geometry.
Conditional on explored geometry and all independent Liouville input the appropriate fresh geometric experiments in cut disks then use the CPI laws just described; in particular their conditional law on each given cut surface only needs the abstract surface/target, permitting iteration under the path disintegrations. \citep[Proposition~5.3]{MSWTrunk}

\proofstage{Proxy covariance for prescribed targets}\phantomsection\label{cut:prescribed-target-covariance}

For use of \eqref{eq:original-11} directly and recursively in these critical laws, let us clarify the random-map metric comparison when an actual target is at a \emph{prescribed} relative length from the typical source. (The laws/contact rules above at such starts used only good approximating subcritical splits, not this critical comparison.) Couple countably many extra typical targets by the ordinary target-consistent bSLE/GFF construction given the independent parent disk.
At every compact initial part strictly before terminal in the prescribed branch there remains a neighborhood of its original target free of the trace (in particular by the contact rule); all targets in a sufficiently small open arc remain unseparated there.
Thus sufficiently close members of the typical list, dense a.s., follow the very same full exploration there, with common clock by positive jump counts and the same cut and post domains.
The immediate parent-to-jump/post comparisons for those countable branches hold almost surely by the earlier fixed-geometry and Devlin argument, since they are ordinary independently targeted in twice-typical parent charts and now have the critical cut/post disk field marginals.
On a common portion one can thus use the same comparisons for the prescription.
If additionally changing standard parent or child charts, these differ by conformal maps on the \emph{fixed} standard open domains; all charts used for the chain estimates qualify there in their own marginal laws (fresh typical mark with the source/tip conditional on the specified perimeter data).
Devlin with the proxy charge and the continuous Weyl log-derivative correction, or the local LFPP comparison above for that correction between qualifying marginal fields, therefore identifies the length elements also across this change.
This proves the needed comparison within the prescribed-start marginal, hence also when this step is used at countably many nodes conditional on abstract previous perimeter data giving that marginal.
Multiplication of the successive local derivative weights uses only actual open interiors.
In their product the upper bound in \eqref{eq:original-11} comes from the \emph{total} composition map into a bounded original chart, without per-generation distortion constants. \citep{MSWTrunk,Devlin}

\end{proof}

\Needspace{5\baselineskip}
\begin{lemma}[A complete shrinking tiling]\label{cut:tiling}
Recursively explore every jump interior from its pinch toward a fresh
independent boundary-length target.  The resulting critical disk tiling has
mesh tending to zero almost surely.  Every macroscopic CLE loop occurs as a
positive jump disk, and area is the sum of the jump-interior areas at every
node.  For some \(p<2\), the sum of the \(p\)-th powers of the perimeters in
one negative-only generation has a strict mean contraction after averaging
over the initial uniform split.
\end{lemma}
\begin{proof}\leavevmode

Here is one useful full tiling by such branches. In fresh jump disks always use the pinch as first mark and a fresh independent length target, sampling initially typical marks. Recurse in all jump interiors. For negative-only steps there is strict moment contraction in perimeter. At \(S=1\), compute twice the killed generator on \(S^{p-2}\) with the compensated added cost of negative offspring sizes to power \(p\) under the transform. It is
\[
 \sum_{\ell=x,y}
 \left\{{\rm PV}\int_{-\ell}^\infty\frac{(1+u)^{p-2}-1}{u^2}\,du-\ell^{-1}
 +\int_0^\ell u^{p-2}(1-u)^{-2}\,du\right\}.
\]
Independent of the split by differentiating (symmetric endpoint terms), this equals \(-2\) at \(p=2\), hence stays negative just below 2.
After division by \(S^{-2}\) the drift/cost rate is \(-c_p S^{p-1}\).
Stopped compensation on compacts and exhaustion give \(\mathbb E\sum_{\Delta S<0}|\Delta S|^p\le1-c_p\mathbb E\int_0^T S_t^{p-1}dt\); hence averaging over the uniform split gives a factor strictly less than 1.
Positive offspring power sum also has finite mean by its compensator proportional to \(\int S^{p-1}\).
Full-path offspring independence as used in iteration follows from the marked rule and geometric domain Markov (future in the post disk independent of past offspring given the post surface and length data); exhaust by compact times increasing to terminal.

Using \eqref{eq:original-11} with \(\lambda'\xi_0=p\) at criticality gives summable upper-chain probability bounds for macroscopic diameters in the negative-only generations by this contraction (apply it in the critical marginal laws and compose comparisons to the same ancestor parent chart).
Thus those generations shrink almost surely uniformly in Euclidean diameter, by Lemma~\ref{proxy:lower-cost} and its almost-sure positive proxy threshold for fixed macroscopic size.
Every infinite chain with infinitely many positive steps likewise shrinks by nested CLE local finiteness in the closed disk; here one uses the actual independent CLE critical coupling just described.
Countably many negative-only subtree starts and the shrinking of siblings at every fixed node by \eqref{eq:original-11} and power sums now imply uniform mesh shrinkage down the full tiling (otherwise follow an infinite large-diameter chain).
All true macroscopic loops occur as positive disks, since a not-yet-hit CLE loop has to survive intact in remaining complementary domains.
Area adds exactly over jump interiors at each cutting: off the outer boundary the complementary skeleton has zero area, by Lebesgue nullity of the independent critical branch trace and interior diffuse LQG that charges no independent Lebesgue-null set (as in Section~\ref{sec:continuum}).
Use typical targets or their common pre-fork explorations.
Disk boundaries themselves carry no area.
Thus the joint domain tiling and its volume labels refer throughout to an actual critical field disk with independent signed CLE, without needing determinism of that quantum surface from any encoding processes.

\end{proof}

\subsection{The loop Palm law of the prescribed sphere}\label{cut:palm-section}

Here we mark a loop by \emph{counting measure} among all nesting depths. When calling one side blue, choose fairly between the two sides in a one-loop Palm law (globally one may color alternately across loops with a uniform parity choice). This side designation is not the sign in the bSLE coupling.

\begin{proposition}[Critical sphere cuts]\label{cut:palm}
For the unit-area critical sphere of Section~\ref{sec:continuum} with an
independent nested \(\CLE_4\), the counting intensity of a loop, a fair choice
of side, its intrinsic perimeter \(s\), and that side's area
\(v\) has density
\[
 \frac1{2\pi^2}s^{-5}d_s(v)d_s(1-v)\,ds\,dv,
 \qquad s>0,\quad 0<v<1.
\]
The corresponding actual cut surfaces are the critical disk pair with these
areas and perimeter, with independent uniform relative boundary phase and
independent nested CLEs in the two interiors.  Matching is by the intrinsic
boundary-length labels.  The identity also holds, after integrating scale,
as a measure-class identity for arbitrary total area.  No uniqueness theorem
for critical conformal welding is required.
\end{proposition}
\begin{proof}\leavevmode

\proofstage{The strictly subcritical identity}

We use the Ang--Holden--Sun SLE loop welding theorem together with the SLE loop measure/whole-plane CLE intensity identification (see Ang--Cai--Sun--Wu): in the simple CLE range \(8/3<\kappa=\gamma^2<4\) the unmarked quantum sphere measure times the independent canonical SLE loop measure (up to constant the expected counting intensity of Kemppainen--Werner whole-plane nested CLE) corresponds on cutting to the uniform-boundary-rotation weld with law proportional to \(\int s\,ds\,{\rm QD}(s)\otimes{\rm QD}(s)\), where \({\rm QD}(s)\,ds\) disintegrates unmarked ordinary \(\gamma\)-quantum disk measure.
We only need the cut and actual boundary identification direction.
The DMS twice-marked disk intensity was used above, so the unmarked density has mass a constant times \(s^{-2-r}\); likewise unmarking the ordinary twice-marked sphere reweights by inverse area squared, irrelevant under area disintegration.
These public inputs are used strictly before the endpoint.
Constant shifts and deterministic measure units preserve this identity up to a scalar, so at \emph{unit} area of the ordinary sphere in \eqref{eq:original-1} it yields the one-loop split intensity density \(C_r s^{-2r-3}d_s^r(v)d_s^r(1-v)\,ds\,dv\), where \(v\) is blue volume.
Conditional surfaces are the two disks with these volumes and perimeter and independent uniform phase alignment (equivalently match independent typical boundary marks), as abstract matched surfaces, not independently positioned in the sphere.
Conditioning here and below at unit total area can be done by continuous area scaling via field constants in the infinite measure identity.
Conditional on the marked loop in the nested intensity the remaining geometric CLEs are independent ordinary nested CLEs on both sides by the whole-plane nested CLE interior Markov and inversion/Möbius invariance (these also give independence from the field before conformal pullback). \citep[Theorem~1.1]{AHSwelding}\citep[Theorem~7.1 and Proposition~7.2]{ACSWloops}

Next restrict to loops separating the two cylinder ends. This adds the factor \(2v(1-v)\) on the intensity side, with the ends then area-marked in the respective disks. Write
\[
 m_r=C_r\,2v(1-v)s^{-2r-3}d_s^r(v)d_s^r(1-v)\,ds\,dv .
\]

\proofstage{Stationary chain intensity and loop spans}

On the cylinder side of \eqref{eq:original-1} we have an independent whole-plane CLE, and now just count the chain enclosing the origin of the plane chart \(z=e^u\).
Write \(t\) for the log conformal radius at 0 of such an interior.
The Palm density in \(t\) is \(I_\kappa=1/\mathbb E\tau_\kappa\), with \(\tau_\kappa\) the disk outermost log conformal-radius drop at a fixed point.
Indeed the intervals to the next smaller loop partition the real line; their Palm drop law is the disk CLE law by nesting.
Translation invariance and this mass transport give the formula (finiteness follows as well by charging an initial bounded interval of each gap).
Here we use the whole-plane nested Markov construction of Kemppainen--Werner \citep{KWnested}, in particular nesting toward 0 after a chain loop is independent given that loop and the exterior, also for the translation Palm intensity (equivalently stop after descending across deterministic radius bounds and count).
By the Schramm--Sheffield--Wilson conformal-radius theorem \citep[Theorem~1]{SSWradii} these drops have uniform exponential tails near \(\kappa=4\), and \(I_\kappa\to I_4=1/\pi^2\) (the limiting transform at argument \(\lambda\) near zero is \(1/\cos(\pi\sqrt{2\lambda})\)).

In cylinder horizontal coordinate the whole boundary of the typical chain loop lies in \([t-w,t+w]\) with \(w\) of stationary law having some uniform positive exponential moment.
One flank follows by the quarter theorem.
For the other, the normalized steps from containing parent interiors to the child containing 0 (interior charts with derivative positive at 0) are iid ordinary disk CLE steps backwards for any finite number of generations under translation Palm.
Indeed shift the count from a loop to its descendant by that number of generations and integrate \(t\) against Lebesgue, using the forward nesting property.
Each step has probability bounded below to lie inside a fixed compact subdisk of its parent chart: use simple loop-soup monotonicity up to 4 at the center.
Go back to the first such good step; its image size in the plane is bounded by a constant times that parent's conformal radius by distortion.
The geometric generation bound and the exponential moment of each log-radius increment (use a small exponential tilt/Hölder) prove the assertion.
Also the law of the normalized interior map at fixed \(t\) converges on interior compacta to the corresponding Palm map at \(\kappa=4\).
Indeed each normalized step has kernel convergence in law by the increasing soup and the conformal-radius theorem.
Compose finitely many; on compacta the normalized map from a parent many generations back contributes an asymptotically negligible nonlinear part by the distortion theorem and the product of the radii ratios tending to zero uniformly in probability.
This characterizes the limiting map and loop interior, independently of the fields.

\proofstage{Determining the normalization constant}

First let \(r>1\) fixed close to 1 and compare expected counts with both volumes \(\ge\epsilon\). Integrating \(s\) explicitly gives leading order
\[
 2 C_r\Gamma(r-1)\Gamma(r)^{-2}\log(1/\epsilon).
\]
In the cylinder, normalized area exponentially far in either end has almost logarithmic rate \(-\gamma a|t|\) there: more precisely the logarithm of the small-side area of the chain loop at \(t\), divided by \(|t|\), tends in probability to \(-\gamma a\) as \(|t|\to\infty\) under the tilted Palm law.
Indeed sandwich by radial tails at \(t\pm w\), use \(R_a(t)/|t|\to a\), and the uniform finite small moment on unit slabs of the lateral mass from Section~\ref{sec:continuum} (and its stationary positive law for lower bounding by one slab). \(w\) is tight.
Moreover on rescaling \(t=z_0\log(1/\epsilon)\), the count probabilities are integrably dominated for large \(|z_0|\): except for exponentially small probability of a huge span \(w\), use the area in the tail beyond \(|t|/2\); its unnormalized small moment decays exponentially by the nonzero radial drift and the slab estimates.
Normalization and the \(L_\gamma^{2a/\gamma}\) tilt preserve an exponential bound by truncation or Hölder; \(L_\gamma\) has the requisite small positive moments with slack near criticality as shown earlier, and its inverse has uniform positive moments (any fixed order by the same packing/inverse argument on one interior slab).
This gives the other leading coefficient \(2 I_\kappa/(\gamma a)\).
Hence
\[
 C_r=r\Gamma(r)I_\kappa/2.
\]

\proofstage{No loss of perimeter at zero}

We check carefully that the boundary length \(s\) is not lost at zero for loops at finite \(t\) as \(\kappa\uparrow4\).
For a translation Palm loop about \(t\) before the sphere field tilt, let \(B_\gamma\) denote its length in the cylinder \emph{lateral field alone}, in the same units as \(s\).
This length is well defined Palm-a.s. at each parameter under consideration by the disk cut rule, undoing the continuous radial and normalization factors.
It is a function of the lateral field and the loop alone, since one can read the seam from the disk interior field pullback via the boundary chaos construction, continuously corrected by these factors.
Its law jointly with \(w\) is stationary in \(t\) and independent of \(R_a\) before tilting.
We have bounds
\[
 B_\gamma L_\gamma^{-1/2} e^{-(\gamma/2)\sup_{[t-w,t+w]}R_a}\ \le\ s\ \le\
 B_\gamma L_\gamma^{-1/2} e^{-(\gamma/2)\inf_{[t-w,t+w]}R_a}.
\]
For fixed \(p>0\), \(\int s^p m_r\to\int s^p m_1<\infty\) by the explicit formulas.
Integrating the lower estimate on \(|t|\le1\) and applying Hölder thus gives \(\sup\mathbb E B_\gamma^{p_0}<\infty\) for some small fixed \(p_0>0\).
In detail use \(B_\gamma\le s L_\gamma^{1/2}\exp(\gamma\sup_{[-1-w,1+w]}R_a/2)\), allocate \(s^{2p_0}L_\gamma^{2a/\gamma}\) to the cut-law integral, and bound the remaining small powers of \(L_\gamma^{\pm1}\) and exponential supremum by the cited moments and the span exponential tail independent of the Gaussian radial input.

Along a subsequence with \(B_\gamma\Rightarrow B\), use the exact identity
\[
 p^2\int s^p\,dm_r/I_\kappa
 =\int_{\mathbb R}\mathbb E[L_\gamma^{2a/\gamma}s^p]/\mathbb E L_\gamma^{2a/\gamma}\,dT_0,\qquad t=T_0/p^2
\]
(where \(dm_r\) denotes the above density).
The integrands as \(\gamma\uparrow2\) are dominated by \(C(1+|T_0|)^{-3/(2q)}\), for some \(1<q<3/2\), uniformly for small enough \(p>0\).
Indeed allocate Hölder power \(q\) to \(\exp(-p\gamma\inf R_a/2)\), with the other powers accommodated by the above bounds.
Its expectation to power \(q\) is \(O((1+p^2|t|)^{-3/2})\) using independence of \(w\) and the three-dimensional Gaussian density bound at \(|t|-w\) when \(2w\le |t|\), integrating the radial exponential and allowing oscillation over \(2w\) (small-\(p\) exponential cost integrable over \(w\)).
The other case is negligible for \(p^2|t|\ge1\).
Thus in particular finite positive-\(s\) windows have uniformly negligible contribution from \(|t|\to\infty\).
At fixed \(T_0\ne0\), in the double limit first \(\gamma\uparrow2\), then \(p\downarrow0\), the upper estimate integrand has limsup at most
\[
 {\mathbb P}(B>0)\,\mathbb E e^{-|W_{|T_0|}|}.
\]
Indeed one drops \(L_\gamma^{2a/\gamma-p/2}\) with an error tending to zero by its uniform moments, and replaces the infimum by \(R_a(t)\) with vanishing error by tightness of \(w\) and Brownian increments; independence and \(\mathbb E B_\gamma^p\to\mathbb E B^p\) apply.
The last Brownian factor integrates (over both half-lines) to 4 by the three-dimensional Green function.
The same constant follows from the explicit density:
\[
 \frac{p^2}{I_4}\int s^p\,dm_1
 =p^2\int_0^1\int_0^\infty
       \frac{s^{p-1}e^{-s^2/[v(1-v)]}}{v(1-v)}\,ds\,dv
 \longrightarrow4.
\]
Hence \(B>0\) almost surely.
In particular in bounded \(t\)-windows the critical approach retains tight strictly positive \(s\), under the sphere tilt too.

\proofstage{Convergence of the abstract matched pair}

Here are details passing the Palm cut rule itself.
On compact positive \(s\)-windows its abstract pair side converges to the formula \(m_1\) with the actual limiting critical disk laws and uniform phase weld alignment.
This does not require pointwise convergence of area-conditional disk laws: disintegrate the homogeneous perimeter integral by starting with independent unit disks and rescaling to \(s=(v_1+v_2)^{-1/2}\) with \(v_i\) their unit-perimeter volumes, reweighting by the power Jacobian from \(s^{-2r-3}\,ds\) and by separation.
This reweighting is bounded and continuous on compact \(s\)-windows up to movable cutoff boundaries.
Endpoints become area samples on opposite sides.
Disk shapes/areas, added chart marks, and phase matching of boundaries by length all converge there.
On the sphere side in each bounded \(t\)-window one has the joint convergence of field/area and the normalized \emph{interior} Palm map just proved, independently.
There is no loss escaping \(t\) or \(s\) in matching these two statements by the tightness/tail bounds above.

\proofstage{Convergence on the actual closed boundaries}

We spell out why the matching transfers on the actual closed boundaries (kernel convergence alone might leave a spurious shear).
Restrict to bounded \(|t|\), \(w\), and \(s\) bounded above and below.
Both disk maps from their canonical charts, in the \(z\) or \(1/z\) chart toward the respective separated end, are bounded univalent maps.
Along the boundary, whose images are in a compact annulus about that end, \eqref{eq:original-11} gives uniform equicontinuity in probability.
Indeed apply the upper chains in the two independent standard boundary charts of the pair, whose laws on these tests before placement are merely boundedly weighted as just explained.
For the proxy lower test only paths in a slightly larger compact parent annulus matter; costs there are uniformly positive for macroscopic motion by Section~\ref{sec:continuum}'s cylinder representation with tight continuous corrections and the fixed proxy theorem.
Even if an upper path extends toward the end itself, it already pays for a macroscopic traversal within the annulus.
We compare metrics only away from the two ends.

For precision about this comparison under re-embedding, first pull back the parent proxy by an interior-normalized map sending the separated end to the origin of the disk in inverse coordinates, with deterministic angular convention from the loop.
Conditional on the Palm loop, the parent cylinder law before \(s\) cutoff is independent, a whole-plane GFF plus continuous up to change of law throughout the open punctured region as needed (via the average-process correction as in Section~\ref{sec:continuum} and conformal GFF comparison).
Thus fixed-map covariance applies on the punctured disk.
The actual maps from twice-typical strip coordinates may also first be recentered at the separated bulk mark, now differing on the fixed punctured disk just by rotation, covered by Devlin.
Their transformed field laws qualify on the punctured disk in this comparison: indeed in a standard child strip the bulk mark location \(z'\) is an area pick.
At these strictly subcritical parameters, conditioning on its location, the lateral field law after subtracting the insertion singularity \(\gamma\log|\,\cdot-z'|^{-1}\) (cut off smoothly locally as desired) is absolutely continuous up to continuous corrections of a whole-plane GFF in the open strip.
This follows by the subcritical Gaussian Palm/Cameron--Martin insertion formula on interior compacta (tilt by a mollified area weight and pass using local subcritical \(L^1\) for the Gaussian field); the covariance shift for \(P^\perp\) has precisely that singularity with continuous remainder, and radial, normalization, total-area and Jacobian weights in the disk pair only change the continuous part or absolute continuity.
Thus, after recentering using a fixed convention given \(z'\), the charge field on the punctured disk, also with the proxy log-derivative correction, qualifies by domain comparison.
Finally one can undo recentering at \(z'\) and compare to the strip metric in \eqref{eq:original-11}: on the whole strip both the actual unconditional field and the singularity-subtracted one qualify by the disk law and the just stated Palm calculation respectively.
Away from the possibly random \(z'\) their LFPP local length elements differ just by the Weyl factor, by marginal convergence as before, while for fixed \(z'\) the adjusted field obeys ordinary covariance.
Chaining these comparisons proves the derivative-weight bound for \eqref{eq:original-11} on all needed path parts before approaching the marks. \citep{Devlin,AHSwelding}

Consequently both maps are tight on closures (maximum principle), and cannot collapse by area and parent diffuseness.
Interior field pullback convergence identifies the actual cut disks (off marks suffices for masses; distribution laws have no point-supported ambiguity, e.g. by local \(H^{-1}\)).
The interior map based at the left end covers exactly the Palm CLE\(_4\) loop interior by its previously proved kernel convergence and bounded-span/Jordan limit.
More explicitly the limiting conformal map from the child sends its stable area mark to the end, hence has this same kernel.
Both boundary images coincide by uniform matching in the cut pair; the exterior map, taking the other end in its interior, therefore covers the complementary Jordan disk.
The limiting parametrizations inherit exactly the cyclic length matching since the disk boundary measures converge to diffuse measures of full support.
Thus the abstract matched critical pair law is attained by actual sphere cuts.
Their critical lengths are functions of the cut field itself via disk pullback, so there is no residual choice along subsequences.
Exhausting the windows in the two descriptions proves the critical Palm equality, with
\[
 C_1=1/(2\pi^2),
\]
for actual cuts separating the two ends of \eqref{eq:original-2}.

\proofstage{Removing the separation marks}
Add the independent nested CLEs on both sides by the geometric conditional law if desired.
Removing separation just divides by \(2v(1-v)\): by Section~\ref{sec:continuum} the ends on the unmarked sphere are area samples, each loop has strictly positive volumes on both sides, and the pullback identification/lengths once valid on separation are intrinsic to field and loop.
For clarity about the CLE variables here, under counting Palm on loops separating given deterministic points the ordinary inside nesting law holds by summing the nested Markov assertion (one can start descending across deterministic containing-radius cutoffs and then exhaust).
By Möbius invariance the same holds from the other side viewed toward the second point; thus conditional on the chosen loop the two side CLEs indeed have the independent nested-domain laws.
Independent cylinder/field sampling retains this rule and end separation exhausts intensity upon area re-sampling as stated; the orientation/blue designation is fair.
Other normalization marks for the embedded sphere may be sampled subsequently by intrinsic area.
This gives the precise all-loop intensity version matching \eqref{eq:original-3}.
By integrating over scale it gives also the measure class identity between the arbitrary-total-area sphere Palm law and two equal-perimeter welded ordinary critical disks with random perimeter of positive density.
No uniqueness of conformal welding at criticality is invoked.

\end{proof}

\subsection{Exact peeling operations and the Cauchy limit}\label{cut:peeling}

We first construct a coupling at the level of the cuttings.
It gives positions up to small-mesh topological approximations, \emph{not yet} conformal positions, and uses no convergence theorem for inventory words.
Write \(T_p\) for the probability disk in \eqref{eq:original-5} at \(X\), with a specified boundary color. A "slot" is an edge occurrence, in particular not a distinct-edge count.
Extra fair independent orientations can be assigned to rings/loops.
Opposite choices below at a discovered ring are equiprobable, independently before exposure; one can make them by these orientations (use the traversal with that orientation in the positive aperture convention of \(Z\)).
In particular all branches with a common source in a disk use the \emph{same} choices, not new choices each time they use a common prefix.
The central ring of a pair construction can separately be marked by a fair choice.
Orientations of different loops are assigned independently conditional on the underlying discrete sphere.
They differ from primal/dual color labels.

On a fresh polygon with active and target boundary edges distinct as slots, keep two ordinary lists of slots, one on each side from the target toward the tip, of sizes \(L,R\); \(S=L+R+2\).
Vertex occurrences along each, including the appropriate target and active endpoints, have heights \(0,\ldots,L\) or \(0,\ldots,R\).
With zero list size the two endpoints on that list are the same occurrence here.
Fix clockwise order in using the side convention, both discretely and for corresponding \(Z\).
Here is the step until too close to absorption:
\begin{itemize}
\item If peeling the gasket deletes the active edge as a bridge with its other copy on (say) the \(L\)-side, let \(j\) ordinary slots be in between on that side.
Keep the target disk.
Take as new active edge the nearest edge on that same side beyond the duplicate, so \(L'=L-j-2\), \(R'=R\).
This instruction is for \(j+2\le L\).
The detached \(T_j\) has its boundary based at the other end of the removed bridge (the old \(L\)-endpoint); either consistent choice of a slot there for rooting, by occurrences, gives its fresh marked law.
Boundary heights on it before deletion are \(L-j,\ldots,L\).
The pre-height \(L-j-1\) is at the \emph{retained} end of the bridge, also the unchanged \(R\)-tip endpoint.
The end at \(L'=L-j-2\) is adjacent to this by the new active edge.
The three roles can be read from the pre-step list as
\[
 \underbrace{0,\ldots,L-j-2}_{\text{new ordinary list}},\quad
 \underbrace{L-j-1}_{\substack{\text{retained bridge end}\\\text{and unchanged opposite tip}}},\quad
 \underbrace{L-j,\ldots,L}_{\text{detached disk boundary}}.
\]
The bridge joins the last two groups; the new active edge joins the first
group's endpoint to the retained bridge end.  Thus a surviving occurrence
at height $L-j-1$ can switch to the unchanged opposite tip.  These exact
incidence shifts will be used in Lemma~\ref{cut:incidence-consistency}.
\item At a face step of degree \(k\ge2\), replace the active edge in the frontier by the other \(k-1\) slots of that face, choosing as new active the one at the unchanged-side end, so (for a discovery on the \(L\)-side) \(L'=L+k-2\) and \(R'=R\).
This is before any intrinsic pinches of boundary labels, which do not change the list instruction.
The added interval follows the circumference from the deleted tip to the new tip, omitting boundedly many slots at the ends as indicated.
The ring's inner disk \(T_l\) is taken based at the matching pinch location in its own list (nearest base according to a fixed occurrence convention), with reverse gluing orientation.
For \(k=1\), delete the edge and use the adjacent old slot on the chosen side as new active (\(-1\) on that side); the old endpoints coincide.
\item If these changes would hit or include the target slot as active or paired edge or cannot continue by this rule, stop; one may leave an entire last portion unexamined. Nothing in a compact-path assertion runs through that exception. Also stop for coincident initial tip and target slot. For several targets, at a split continue in the polygon of each target by the same rule for its path (in particular from the pinch, using the side-dependent adjacent edge in that polygon).
\end{itemize}

\proofstage{Fresh fillings and vertex incidence}

The ring/bridge rules are inside \eqref{eq:original-5}, thus every as-yet-unopened polygon filling has fresh Boltzmann law, rooted by its slots, conditionally independently given the skeleton information (including exposed choices and lengths).
This extends to full branch paths for the jump fillings by successive conditioning: further trajectory choices along the path to its target do not look inside the detached fillings.
The remaining target polygon itself has this rule when stopping by trajectory data.
Every one of these polygons has its vertex set injecting into the assembled map if the original polygon did.
Indeed a bridge joins \emph{disjoint vertex sets} of its two disks; in the face case one attaches a face along frontier slots and adds an edge without identifying two old vertices.
The opposite-color ring boundary is an independent list filled by \(T_l\); it carries no extra identifications of its vertices by attachment on the outer side.
In the sphere joined by a ring both initial vertex sets likewise inject.
Boundary pinches in all these assertions are intrinsic to the fillings, not identifications of distinct vertices by an unknown exterior equivalence.
These facts follow also by doing each inverse operation before attaching all other rings.
Sharing with pieces external to a filling is possible only via its boundary vertices (the single vertex for perimeter zero).

\begin{lemma}[The raw walk and localized peeling convergence]\label{cut:raw-walk}
The nonabsorbed peeling path is a raw iid walk, restricted to admissible
side updates and weighted by \(f_{S_{\rm end}}/f_{S_{\rm start}}\).
With slots divided by \(P\) and time unit
\(m_P=(\pi^2/2)P/\log P\), it converges on every compact surviving path test
to the two-sided critical law of Proposition~\ref{cut:critical-cutting}.
This convergence includes the signs, sides, perimeters, and ordered aperture
labels of every fixed matched jump, and can be iterated through fresh jump
fillings.
\end{lemma}
\begin{proof}\leavevmode

On nonabsorbed paths the increment law is exactly iid with a likelihood ratio \(f_{S_{\rm end}}/f_{S_{\rm start}}\), requiring admissibility of the side updates, with iid weights on \emph{each} side
\begin{equation}\label{eq:original-12}
w_{-j-2}=c_*^{-2} f_j,\quad
 w_{k-2}=c_*^{-2}\widetilde g_k/2 \quad(j\ge0,\ k\ge1),
 \qquad \widetilde g_k=c_*^k g_k(X).
\end{equation}
The coordinate not indicated by the side stays fixed.
These are normalized (\(2\sum w=1\)): divide \eqref{eq:original-5} by \(F_p\) and take \(p\to\infty\), splitting the first convolution at its midpoint (either convention for a midpoint term, of vanishing effect) and using dominated convergence by the two-sided estimates on \(f\).
In the face case \(l\) and the interleaving have the tilted uniform-ring law in \eqref{eq:original-5}, i.e. a negative-binomial allocation tilted by \(f_l\) with mean parameter \(k\) before tilt.
In particular for \(k\to\infty\), \(l/k\to1\) and the order from a fixed outer base with the fixed starting-link convention identifies equal \emph{relative} length positions uniformly with error tending to zero in probability.
Indeed exponential errors at any fixed fractional tolerance for totals and partial sums of the geometric variables beat the polynomial reciprocal normalizer.
The equivalent order in the two-sided central ring, given \(k,l\to\infty,\ l/k\to1\), follows by sampling without replacement.

Take length unit \(P\to\infty\), measuring slots divided by \(P\), and time unit \(m_P=(\pi^2/2)P/\log P\).
Under the raw weights in \eqref{eq:original-12} we get \(J_1\) convergence on compacts to two independent symmetric Cauchy processes in exactly the units above (densities \((1/2)dx/|x|^2\) each).
Indeed the signed truncated first moment at symmetric cutoff is bounded, by \(\sum k|\widetilde g_k-2f_k|<\infty\) as proved above.
Jumps outside a neighborhood of zero converge by the iid Poisson limit; smaller compensated sums vanish uniformly in probability in the small-cutoff limit by the truncated second-moment bound, and the residual bounded drifts are negligible.
Thus from any convergent positive pair of initial ordinary lengths, on survival through any compact piece the discrete limit is given by the likelihood \((S_0/S_t)^2\) and the two-coordinate positivity restriction, i.e. the critical disk rule of the preceding sections (\(S\) here in scaled units).
More precisely one can test horizons with both paths bounded and bounded away from zero, by path tests with null boundaries; these restrictions are convergent finite measures, not necessarily laws of the whole path.
On such restrictions the exact finite ratio is bounded and convergent by \eqref{eq:original-12} and the \(f_p\) equivalent.
All \(0,\pm1,\pm2\) slot conventions above matter exactly on the lattice but do not move macroscopic jump times or aperture labels in this convergence.

Here and below one can couple, along further subsequences, with a true critical path on its entire exhaustion by compacts before \(T\), matching every fixed positive compact time segment in \(J_1\) eventually.
Nothing is required of the terminal discrete path duration.
For detail, take horizon-and-range tests, with continuity endpoints, inset them if necessary and increase a countable family covering every compact surviving initial segment in the critical law.
Convergence on such tests including their path coordinates holds also on finite intersections (use the longest surviving horizon) and on differences by subtraction.
The critical law is probability, staying strictly within positive finite length ranges on compacts of \([0,T)\), and extending to side sum zero at finite \(T\) as above.
Thus the localized descriptions exhaust it; elementary weak-convergence coupling on increasing finite surviving tests and diagonalization give the assertion (any leftovers of mass tending to zero can be arbitrary).
Time alignment matches individual jumps, their signs and sides.
Child negative perimeters are \(j/P\), positive inner perimeters \(l/P\); both converge to the jump magnitudes.
The matching extends to boundary label orders in the indicated direction, including the inner bases across the face ring.
It extends recursively to every fixed matched jump filling by the product Markov rule and the same convergence for arbitrary converging starts.
One convenient primary experiment initially chooses an independent uniform target, and does so afresh in every jump interior based at the pinch.
Degenerate lattice marks can just stop that polygon.
We always compare orientations from a common cyclic traversal, so the aperture order of the persistent ordinary list corresponds to the continuum order from the target back to the active tip (and on attached disks to the order along that arc).
This is the clockwise/counterclockwise choice of the discovered-loop traversal in the exploration convention; interchanging left and right everywhere just reflects both descriptions.

\end{proof}

\begin{proposition}[Compatible target routes]\label{cut:marked-targets}
For any finite compatible family of transmitted or independently sampled
interior target labels from a common source, the localized path and offspring
comparisons hold jointly, for arbitrary converging positive initial
perimeters and splits.  The same holds recursively in fresh matched jump
disks, with the prescribed pinch source or a single fresh common source;
transmitted targets are chosen without inspecting the fresh filling.
Countably many such observations are obtained by increasing finite tests.
Only compatible routes use common prefixes; independence is asserted for
fresh components under the marked conditional laws.
\end{proposition}
\begin{proof}\leavevmode

A useful extension is to several prescribed or randomly sampled target labels from one active source.
Before the loss of a target label along a reference path the prefix containing it is unchanged.
For a distinct target in the open ordinary arcs initially at strictly interior height \(u\), almost surely the first descent there below \(u\) is a strict negative jump with pre-value above \(u\), and until then the length is uniformly strictly above \(u\).
This uses the zero Lebesgue measure of the closed descending Cauchy ladder height range, its scaling (hence polarity of any fixed positive height difference from the start), localization, and terminal limit zero.
Random labels here can, for example, be conditionally independent given starting lengths, with positions converging in those units.
Thus each fork of finitely many targets is in a compact-surviving piece, with strict slot transmission into a recognized jump.
The prefix is simultaneously a prefix from the source toward all those swallowed targets, until separation.
At the separating bridge, continuation toward them starts in its jump disk; the new active slot at the pinch is prescribed by the rule, the same one for those strictly inside this detached list.
Their lengths there converge to the transmitted labels.
Unswallowed targets are treated along the continuation.
The continuum common-prefix and fresh component laws follow by the branching bSLE/GFF coupling and the strict stack rule proved above; unclocked branches use independent exploration laws in newly split domains (conditional also on the field supplying the quantum endpoints and the past, which is valid by exploration Markov with this initially independent input).
The clock on common portions agrees, e.g. by jump counts.
Fresh abstract disks with perimeter labels are used by the marked rule at the split; any imposed transmitted targets, when also conditioning on past path/split data, do not look inside them.
This proves joint convergence of finite compatible routes by induction, including their offshoots.
One can continue similarly on compact pieces ending before a fork, or select finitely many further paths recursively in matched raw jump disks from their prescribed pinch or a \emph{single} freshly sampled source there, provided the rules are consistent (unexplored disk when resourced).
We use countably many observations of this kind by increasing finite tests.
No assertion that paths from conflicting sources using the same map are conditionally independent is intended.

\end{proof}

\subsection{Boundary pinches and consistency of all occurrences}\label{cut:occurrences}

We next prove the spatial validity of these couplings. We give details because convergence of slot processes with non-simple discrete boundaries need not by itself give a uniform projection of all flags.

\begin{lemma}[No macroscopic pinch on a fresh boundary]\label{cut:fresh-pinch}
In a fresh Boltzmann disk \(T_p\), the maximum cyclic separation of two
boundary occurrences of the same vertex, divided by \(p\), tends to zero in
probability.  The estimate transfers to fixed deterministic surviving times
and to a tight number of selected macroscopic jumps on positive compact
path tests.
\end{lemma}
\begin{proof}\leavevmode

In \(T_p\), index the vertices cyclically by boundary occurrences.
The probability that two occurrences at slot difference \(d\in(0,p)\) are equal is \(f_d f_{p-d}/f_p\): cutting at the two corners of this common vertex gives exactly two independently weighted disks when this equality is specified, with ordered roots (cyclic rerooting of an original root is a bijection).
This allows all further pinches inside either disk.
Indeed boundary arcs between the cuts and the maps inside those arcs split at the vertex in its rotation, and conversely wedge the two roots' corners at that vertex.
Loop data do not cross the cuts.
Fix two small disjoint index arcs of macroscopic separation \(I<J\) in a cyclic coordinate based off them.
Expected number of equal pairs in slightly enlarged such arcs is \(O(\log p)\).
If some pair in \(I,J\) is equal, choose the one with smallest \(i\in I\), then largest \(j\in J\) for that \(i\).
Given that \(i,j\) are equal, this selection for a specified pair uses only the \emph{outer} disk of the cut (indices outside \((i,j)\)).
Indeed any contender at \(i'<i\) with \(j'<j\) crosses and hence has the same vertex, giving the exterior contender \(i',j\); the case \(j'\ge j\) is already exterior.
Larger \(j\) for the same \(i\) is exterior as well.
Thus the inside is fresh.
Its conditional expected number of additional equal pairs at \(i+s,j-r\), \(r,s\ge1,\ r+s\le \epsilon p\), is at least \(c\log^2 p\), by summing \(f_{r+s} f_{j-i-r-s}/f_{j-i}\).
Choose \(\epsilon\) small in the enlargements.
The unconditional upper bound thus gives probability \(O(1/\log p)\) of any equality.
Finite coverings show that maximum cyclic separation among identical boundary vertices divided by \(p\) tends to zero in probability.
This is not claimed uniformly over \emph{all times} of a peeling.
It holds at a fixed deterministic scaled surviving time with perimeter in positive compact ranges, by freshness.
It holds for the post polygon and all macroscopic jump fillings at selected large jumps there, also under step-type and jump-size selection, by the factorization.
To use simultaneously over jumps above any positive cutoff on localized paths, their count is tight; one can first mark in stopped order up to a fixed count with posited compact bounds.
All these estimates and the analogous countable rational-time, horizon and matched-node estimates may be imposed in the coupling by further subsequences.

\end{proof}

\proofstage{Continuous prefix labels}\phantomsection\label{cut:prefix-labels}

In the quotient description already proved for the critical disk define
\(B^\alpha_v(h)\), \(0\le h\le Z^\alpha(v)\), at a state \(v\) (a pre- or post-real-time position), by the backward prefix: follow that side of the contour backward, including detours and ultimately the increasing arc on the old boundary, to height \(h\), transmitting above \(h\).
Any such linked choice has the same physical image.
These labels are continuous in \(h\) and along converging state copies (ordered with detours).
Indeed along a converging sequence any subsequential limit of backward choices still has the required linking property by compactness on the height-excursion circle and continuity of heights; all choices coincide in physical image.
At matched jump \emph{post} fronts they are also the actual continuing-side perimeter labels in the limiting disk maps.
For open outer/insertion labels strictly transmitted to that state this follows from the subcritical comparisons and uniform convergence of post maps with their boundary measures; non-strict ties from those aperture interiors are obtained by perturbing the height slightly down within them.
If tied to a real time on following backward, undercut by a large constant times mesh oscillation errors (large enough to give a strict margin, with the perturbation tending to zero), working from a nearby post partition time at that earlier end as in the linkage-to-tip proof above; at endpoint jumps use the post-state/aperture endpoints.
The starting images then approach that tip by the simultaneous short-arc estimates.
This also identifies such labels if keeping an ordinary deterministic-time post map by the later-jump approximation described above, since the later boundary measures in standard coordinates have limiting disk laws jointly with their fields and hence identify the actual boundary length by field determination and covariance.
Pre-jump notation by contrast can just use the contour/link definition.
On the nonjumping side one can just hold the position over the other side's jump.
At the top the image is the trunk \(P_{\rm tr}(t)\), also for either copy at a jump; at height zero it is the target.
Starting labels are the initial boundary points.
Inserted positive apertures agree with \(B\) just post, detached negative apertures with \(B\) just pre.
Two prefix labels at the same height across states staying at or above that height in between on that side have the same image.

\begin{corollary}[Finitely many labels on a fixed frontier]
\label{cut:finite-frontier-labels}
Fix a pre- or post-state $v$ of a critical branch.  For every physical point
$z$, the set
\[
 \{(\alpha,h):\alpha\in\{L,R\},\ 0\leq h\leq Z^\alpha(v),
                   \ B^\alpha_v(h)=z\}
\]
is finite.  This assertion concerns one fixed frontier.  It does not bound
the number of historical contour occurrences in a contact class.
\end{corollary}
\begin{proof}
Within either prime side, a direct stack class has a single height, so it
contributes at most one label to that side of a fixed frontier.  Distinct
heights can be joined only through the folds at a real trunk time.  The
exact contact rule permits only one such time in a physical class; that
time has finitely many pre/post side copies, with at most one nontrivial
linkage attached.  Hence only finitely many heights can occur.  The source,
target, and aperture endpoints add only their prescribed finite copies.
\end{proof}

\begin{lemma}[Consistency of vertex occurrences]\label{cut:incidence-consistency}

Consider one paired parent path.
For arbitrary vertex occurrences at compact-surviving lattice states whose times (with \(J_1\) alignment on compacts), sides and rescaled heights converge to \(v,\alpha,h\) as above, use the proposed location \(B^\alpha_v(h)\).
Pass to subsequences throughout for pre/post alignment.
Two such occurrences of the \emph{same vertex}, even at different times, have identical proposed locations.
Moreover if this vertex survives to the portion beyond every compact initial piece, its location must be the target.
The assertions hold simultaneously as subsequential tests, almost surely on our good coupling subsequences.

\end{lemma}\begin{proof}\leavevmode

\proofstage{Comparisons at a common state}

First compare two occurrences at one common state tending to \(t<T\).
If both heights tend to their respective tip heights there is nothing to prove.
Suppose a label is a positive distance below its tip.
At a regular \(t\), or at a jump approached on the post side, this occurrence (and any other strictly below tip) persists to nearby fixed rational times on the right where both coordinates have moved as little as desired.
At these times the no-pinch bound applies.
Two interior labels must thus have the same proposed value: their circular offsets must converge together (allowing both at height zero on opposite sides).
If the other is instead near-tip, track that occurrence forward while it survives.
Until a first pop below its height, the identical label persists on that side.
At such a pop it either belongs exclusively to the detached disk and cannot survive at all in the target polygon, or it is exactly at new height plus one in the old indexing, hence is on the \emph{other} unchanged top and can switch there (the monogon pop follows this same switch rule).
Since the interior occurrence guarantees survival, one can continue this tracking, always staying within the small oscillation error of one of the two tops until the compared rational time.
Its macroscopic cyclic separation from the interior occurrence (both sides initially bounded below positively) violates no-pinch there.
The same argument applies at time zero.

At a jump approached in pre-states transfer to just pre-jump in the lattice, with oscillations tending to zero.
Again if both labels are near tip their values agree anyway; otherwise at least one interior occurrence guarantees survival to that state, and a possibly switched near-tip occurrence there still limits to the pre-tip on one side.
For a positive macroscopic jump the whole old ordinary lists and endpoints embed as inherited lists in the post polygon with asymptotic lengths \(Z(t-)\); no-pinch there forces identical images of the transferred labels, since they transmit horizontally and the two top conventions have the stated aperture endpoints.
For a negative jump use disjointness of the retained and detached vertex sets.
Within the latter the labels in \([L-j,L]\) (for a jump on \(L\)) have no pinch at macroscopic \emph{cyclic} separation, so limits agree or are the two aperture endpoints, of common image.
In the retained polygon the old heights up to \(L-j-1\) project to the corresponding post lengths, with the last of those heights at the retained bridge end; all horizontal labels through that threshold transmit, the limiting threshold itself mapping to the pinch.
No-pinch in the retained polygon again suffices.
Here \(L,R\) in those interval descriptions are integer pre-state coordinates.
This proves the same-state claim at every possible compact limit.

\proofstage{Completed switches have negligible total duration}

For a longer comparison follow a surviving occurrence by precisely the constant-height-until-switch rule just used, without changing occurrences at other coincidental pinches.
On compact path ranges, after switching to an exact top, the total durations of \emph{completed subsequent switches} before any compared later compact state are \(o(m_P)\), uniformly.
We justify this uniform statement.
Under raw iid steps, the time to first loss of a height starting at an exact top on one side has tail at \(n\) at most \(n^{-1/2+o(1)}\), since the coordinate walk has \(\mathbb P(Y_k\ge0)\to1/2\) starting at 0.
Indeed the Sparre Andersen--Spitzer fluctuation identity \citep[Section~3, Corollary~2, Equation~(3.7)]{SpitzerCombinatorial}, with ties on the nonnegative side, gives the generating function of survival tails as \(\exp\{\sum_{k\ge1}z^k\mathbb P(Y_k\ge0)/k\}\).  For every \(\eta>0\), this is \(O((1-z)^{-1/2-\eta})\) as \(z\uparrow1\); monotonicity of the survival tails and the choice \(z=1-1/n\) give the stated bound.
At first loss, success in switching requires undershoot exactly one, with probability bounded above by \(1-c\) (already the first step can undershoot farther).
Restart this raw rule iid from the opposite coordinate on success. A union bound over either coordinate and \(O(m_P)\) possible starting steps, and the stopping-time product estimates, give with high probability no run of more than \(C\log P\) consecutive successes, nor a run among these with at least three completed durations above \(m_P^{0.9}\): the latter costs at most \((C\log P)^3 m_P^{-1.35+o(1)}\) per start.
Thus apart from at most two longest durations the total completed duration is negligible.
Transfer to localized pieces by the bounded ratios.
Any single completed such duration lying on these compacts is itself uniformly negligible by \(J_1\) and absence of scalar Cauchy chords: otherwise its side heights start at a top and end one below that value, sustaining it all the way before the final pop.
In the limit this is a level linkage between distinct real times including the appropriate copies at endpoint jumps, which was excluded in the quotient proof.
This proves the statement by another countable exhaustion and good subsequence.

\proofstage{Comparisons across times and at terminal}

Now track from an earlier occurrence until the later state.
Survival indeed holds throughout if it appears later.
If there was no pop below it, its proposed earlier value equals the one transmitted to the later state.
Otherwise at the first switch, limiting at some \(r\), the original height has transmitted exactly to a tip copy (error one at the switch), of image \(P_{\rm tr}(r)\).
All subsequent \emph{completed} switches have total time negligible, so the final tracked height arose at a top also limiting to some copy at \(r\), then transmitted without pop through the compared time.
It too therefore has that same image.
Use the same-state assertion between it and any desired later occurrence.
Finally, from a compact-surviving state with positive limiting height one cannot stay in polygons after every compact toward \(T\): without a switch the height contradicts \(S\to0\), hence first switch must itself occur on compact pieces after subsequence.
In a neighborhood of its limit time \(r<T\) both side heights (either copy) are bounded below positively.
On each longer compact horizon the remaining completed switches have negligible duration, hence the tracked height at that horizon still limits to values bounded below by this \emph{same} positive bound.
This again contradicts \(S\to0\).
Height zero maps to the target as claimed.
This argument classifies simultaneous limiting pinches without an estimate summable over every deterministic frontier.

\end{proof}

\subsection{Spatial comparison of the entire flag sphere}\label{cut:spatial-section}

\begin{lemma}[Uniform lifts give nearby homeomorphisms]
\label{cut:homeomorphism-promotion}
Let $X_n$ be topological spheres and $G_n:X_n\to S^2$ continuous maps whose
images become dense in the spherical metric $\rho$.  Suppose that for every
$\epsilon>0$ there are $\delta>0$ and $N$ such that, for $n\geq N$,
\begin{equation}\label{cut:uniform-lifts}
 \rho(G_n(x),G_n(y))<\delta
 \quad\Longrightarrow\quad
 \begin{gathered}
 \text{there is a path from $x$ to $y$ in $X_n$}\\[-2pt]
 \text{whose image lies in $B_\rho(G_n(x),\epsilon)$.}
 \end{gathered}
\end{equation}
Then there are homeomorphisms $H_n:X_n\to S^2$ with
$\sup_{x\in X_n}\rho(H_n(x),G_n(x))\to0$.
\end{lemma}
\begin{proof}
The lifting hypothesis and dense images let us follow any fixed target
path inside an arbitrary open tube, with prescribed lifts near its two
ends: choose a fine chain along the path, choose lifts of its intermediate
points using density, and join successive lifts by
\eqref{cut:uniform-lifts}.  More generally, all inverse images of a compact
path-connected target set lie in one path component of the preimage of
any prescribed open neighborhood, for all sufficiently large $n$.

Fix a fine simple 3-connected triangular grid on the target sphere.
Around each vertex choose nested balls with strict radial buffers; the
largest balls are disjoint and avoid unrelated edges.  In the source,
enclose the entire inverse image of the closed innermost ball in a compact
connected polygonal subsurface $A$, lying in the preimage of a slightly
larger ball $B_r$.  This is possible by the preceding path-component
statement, compactness, and a sufficiently fine source triangulation.
Choose the outer vertex ball $B_R$ with $R>r$.  Apply the same statement to
\[
 K=S^2\setminus B_R,\qquad
 U=S^2\setminus\overline{B_r}.
\]
All inverse images of $K$ join inside $G_n^{-1}(U)$, which avoids $A$.
They therefore lie in one complementary component of $A$.  Fill every other
hole of $A$.  The resulting closed disk still maps into $B_R$ and contains
all innermost-ball preimages in its interior.  The disks for distinct grid
vertices are disjoint.

Choose thin open tubes around the grid edges, intersecting one another
only inside the innermost balls at common endpoints and avoiding unrelated
outer balls.  Lift each edge through its tube with ends in the assigned
source disks.  Polygonal approximation, trimming, and simplification give
simple arcs between these disks, with distinct endpoints on their
boundaries.  Outside the disks these arcs are disjoint: any intersection
of their image tubes would lie in an innermost ball, whose entire preimage
is already inside the corresponding source disk.  Insert disjoint stars
inside the disks to complete an embedding of the grid.  Whitney's
uniqueness theorem for simple 3-connected graphs identifies its faces with
those of the target grid, up to reflection
\citep[Theorem~1.2]{BrinkmannWhitney}.

For each target triangle, choose a slightly larger ball containing its
triangle, edge tubes, and vertex balls.  The image of the corresponding
lifted face boundary lies in a smaller concentric buffer.  All preimages
exterior to the larger ball can be joined while avoiding this boundary,
by the same buffered path argument.  Some such preimages lie on non-face
graph lifts and hence outside the lifted face.  Consequently no point of
the lifted face maps outside the larger ball.  Choose the grid, vertex
balls, tubes, and face buffers with diameters tending to zero.  Extending
the graph correspondence across each face gives homeomorphisms uniformly
close to $G_n$.  A diagonal choice proves the assertion; no control of
parametrizations inside individual filled triangles is required.

\end{proof}

\begin{proposition}[Spatial comparison under coupled cutting data]\label{cut:disk-spatial}

Consider a pair joined by a central ring of two positive asymptotically equal rescaled degrees, matching equal relative slots uniformly.
Take an actual matched continuum pair, a sphere topologically, with the two closed disks glued by the indicated cyclic length alignment.
Use any compatible sphere coordinates, possibly sampled from a continuum conditional law.
In each disk take a recursive routing as above with continuous trunks, Jordan jump disks shrinking at accumulating jump times (including at terminal), exact quotient contacts, and mesh tending to zero down the nested tiling.
This holds in particular for the fresh-target primary experiment.
Suppose the pathwise localized couplings and vertex estimates of the preceding paragraphs hold on this countable recursion including the indicated ring alignment.
At a matched child the path(s) there are treated within the child's own chart; a continuation toward transmitted targets can use the pinch slot specified by the split.
Then one can align the entire flag spheres by orientation-preserving homeomorphisms to the continuum sphere with vanishing flag mesh so that all sites in matched disks project in the limit within their assigned closures, and the contour, tip/frontier and ring observations above compare uniformly as specified.
Full macroscopic loop tours converge with multiplicities and no macroscopic extras if these cutting trees carry the stated CLE, including all remaining loops inside children.
This assertion does not control either graph distance or actual conformal distance in the lattice maps.

\end{proposition}\begin{proof}\leavevmode

\proofstage{Positions from finite truncations}

Here are details.
Use increasing finite truncations sufficiently slowly and further coupling subsequences: compare on surviving compact initial pieces tending to \([0,T)\) at each paired node, with uniform path errors tending to zero after alignment there.
Match finitely many jumps and their ancestry, eventually processing each fixed matched jump strictly before \(T\) in this countable hierarchy; any macroscopic limiting jump on such compacts is separately recognized.
One can now assign positions (possibly several per site) in the continuum sphere.
\begin{itemize}
\item At a compared frontier assign vertex occurrences by \(B\) on their side at the aligned time using scaled heights clamped to the side range, in particular on original lists. States approaching a matched large jump from just before its step use aligning times approaching the pre-copy, those just after approach the post-copy. Thus these assignments have the subsequential interpretation already proved at converging states of each fixed chart.
\item Assign any bridge of such a step to the trunk there.
For a matched positive jump put ring cells at corresponding positions on the matched jump boundary in relative-length order; at an unmatched positive jump assign them to the trunk.
Put all unmatched/unexamined jump interiors (positive or negative) likewise at the trunk of the event.
Here matched but not yet processed child contents can instead always be put arbitrarily within their continuum closure.
Anything past the compact piece in this chart, including a possible unused final polygon, not addressed by a specific compared jump is assigned to the target.
At the central ring use its own boundary alignment.
\item Allow assignments of corners and whole flags through their own Tutte triangles as well as shared-site assignments; one can use a common cell position on each individual triangle through its own piece, up to vanishing error. This still refers to the actual polyhedral \(X_n\), with monochromatic bridges and other polygon boundary slots accessible as paths there.
\end{itemize}

Within a fixed paired parent, at an unmatched jump with alignment limit strictly before \(T\), its scaled outer magnitude tends to zero by \(J_1\) and eventual matching of large jumps; boundary labels there thus limit to a top.
For stable matched events child-parent orders on slots agree asymptotically: for negative compare \([L-j,L]\) on the indicated pre-frontier with the rooted detached boundary, and for positive use the positive post aperture, matching the child across the radial links.
The boundedly many omitted endpoint slots in the ordinary-list convention limit to the pinch.
At a varying sequence of distinct matched child events all positions in descendant closures converge toward their trunk or target limit by the continuum shrinkage.
This applies also if they accumulate at a large jump time, keeping that one stable jump itself separate.

\proofstage{Agreement of competing incidences}

All simultaneous limiting positions through competing incidences of a vertex agree; on flags the position diameter including corner locations tends uniformly to zero.
To see this, argue by subsequential extraction.
Within one fixed parent, compare frontier occurrences of a common vertex by the previous consistency assertion.
It covers also equality to a site still present beyond every compact piece (terminal comparison).
Shared sites of fillings with other pieces go only through their boundary occurrences.
Bridge ends compared on compact times are old tip incidences.
Thus whenever a negative filling or outer side of a ring has to be reconciled with other positions in this parent, use the same comparison at its aperture and tip states.
Inner vertices of a positive ring (only shared with its child) use the relative-order alignment to the child's own boundary, not a parent-frontier equality across that ring.
All competing positions of such a vertex reduce to comparisons within the child and its outer boundary list, or the whole ring/child already limits to the same point if it varies/shrinks or is discarded.
The two disks of the central ring similarly use their own outer boundary comparisons.
More formally descend through stable paired nodes when following two competing addresses.
If they lie in a common child there is nothing to compare before descent.
At their separating level (including a branch with no fixed next child) introduce extra occurrences of the given vertex on any involved stable child's boundary if sharing outside that child is required, and extract their limiting labels.
The parent comparison just given identifies limits there, leaving only comparisons within the children with their own outer occurrences.
Descend again; any unfinished descending entry lies in an arbitrarily small closed cell, by nested mesh.
Thus all comparisons hold in the limit (after separation any further recursion is a comparison to a boundary representative at a specified point, transmitted by equality or contained along with the other position in arbitrarily small cells).
Every actual Tutte face is on a ring or in a discarded portion and its vertices admit its own position up to the stated errors.
This proves uniform incidence consistency by contradiction.

\proofstage{Joining every pair of lifts of one limiting point}

Choose single values and interpolate continuously to maps \(G_n:X_n\to{\mathbb S}^2\) on the flags, e.g. by spherical chord interpolation of the close triples. Their images get dense, by accessible contour positions and nested mesh. We claim also
\begin{equation}\label{eq:original-13}
G_n(x_n),G_n(y_n)\longrightarrow z\quad\Longrightarrow\quad
 x_n,y_n\ \text{can be joined with images tending uniformly to }z .
\end{equation}
Move first to incident vertices.
Any lifts of the same outer or aperture label can be joined \emph{on those lists} along arcs of vanishing relative-label width, including the cyclic endpoint conventions (this refers to actual boundary-slot lifts; to compare other lifts see below).
At a negative pinch one can cross the bridge and new active edge; at a positive pinch likewise the active edges connect the old and new tip without needing to round a macroscopic loop.
Inner-to-outer boundary transfer at a ring uses the matched radial links.
Tip lifts at times tending to one \(t<T\) connect by the consecutive active edges (which share ends at the updated tips), all intermediate steps having projected endpoint limits at \(P_{\rm tr}(t)\) by continuity including the folds.
At the source one uses the first active edge.
At terminal, sites in all sufficiently late portions are inside the same connected polygon at any fixed earlier compact state; the double-limit supremum of its image distance to target vanishes as the state tends to \(T\) by the upper assignments and shrinkage.
This joins tail lifts to target endpoints which themselves have the target edge.

Horizontal linkages transfer by lowering height.
Given a linked pair on a side at height \(h\), for an outgoing open positive-aperture label take the post frontier, for an incoming open negative-aperture label take the pre frontier, and for an outer label take the initial list.
For a real-time occurrence take the copy at its linkage end, using the paired jump if applicable or approximating by nearby states.
Top coordinates there can be joined to the desired tip by the preceding observations also across copies.
With this choice the intervening states sustain height \(h\).
The lattice comparison exactly transmits a lowered height \(\lfloor P h'\rfloor\), \(h'<h\) sufficiently nearby, by \(J_1\) and the step rules on compacts, or height zero without lowering.
Join along each observing frontier back to the desired label, at upper-image error vanishing as \(h'\to h\), by continuity of \(B\).
This permits the lowered height to be outside an aperture if needed at an endpoint; hole-boundary lifts at the exact endpoint itself freely reach the tip by the preceding pinch joins.
Nontrivial linkages occur strictly before \(T\).
Together these transfers and same-label/fold joins implement the exact quotient contact rule.
An ordinary prefix \(B_v(h)\) itself, if observed, can likewise reach its backward linked contour occurrence.

Apply this recursively after piece-address extraction.
In the strict interior of a child only sequences there can limit to the point by the upper comparisons.
If sequences separate at a fixed parent, any involved stable child position must be on its boundary, so introduce a prescribed lift by boundary slots of that child. A varying/shrinking child's interior sites can move to its attachment with all locations still tending to \(z\), by the uniform upper comparison there and connectedness including the bridge or ring.
This leaves outer/aperture or tip/target representations of \(z\) at the separating level, joined by the parent rules just proved by absence of extra contacts.
The central ring case only needs its own length-aligned boundary joins.
Descend to join original lifts to newly prescribed boundary lifts \emph{within} the stable children if necessary, also descending if the two lifts remain in common cells.
This recursion suffices for any finite accuracy: stop along unfinished branches with arbitrarily small image closures and join freely within those connected fillings, using nested mesh and a diagonal choice, with finite joins at each fixed truncation.
This proves \eqref{eq:original-13} along further subsequences and hence in the stated sense.

\proofstage{Promotion to orientation-preserving homeomorphisms}

The sequential assertion \eqref{eq:original-13} implies the uniform
hypothesis \eqref{cut:uniform-lifts}.  Otherwise, for some fixed
$\epsilon>0$ one could choose $n_j\to\infty$ and pairs whose image distances
tend to zero but which have no joining path in the corresponding
$\epsilon$-balls.  Compactness gives a common limiting image $z$, and
\eqref{eq:original-13} contradicts that choice.  Dense images were proved
when constructing $G_n$.  Lemma~\ref{cut:homeomorphism-promotion} therefore
gives homeomorphisms uniformly close to $G_n$.

Their orientation is pinned by the central interface traversal once
around its simple limiting curve, with the specified disk side on its
left in both pictures.  Interior lifts on the two sides approach points
separated from the limiting curve on the corresponding sides.  The
uniform ordered comparison of that traversal preserves its winding, so
the nearby homeomorphisms preserve orientation.  Since the images of
whole flags under $G_n$ already have uniformly vanishing diameter, the
same is true under these homeomorphisms.

\proofstage{Complete loop tours and exclusion of extra loops}

Each fixed positive jump loop, like the central one, traces its matched loop once in relative label order with sign/orientation as observed. Both its outer and inner scaled degrees converge to the same quantum length. Conversely an unnamed loop remains wholly in a yet-unexamined polygon or is itself on a discarded ring; by extraction and shrinkage of varying children, terminal parts and deep stable cells no extras of positive macroscopic diameter occur. This proves the proposition.

\end{proof}

\begin{proposition}[Joint path, volume, and area comparison]\label{cut:joint-volumes}
At perimeter scale \(P=B_n\), normalize triangle counts by \(2n\).
In the unconditioned pair construction, all matched cell volumes converge
jointly with the compatible path data to the areas of the corresponding
critical field disks.  Under the homeomorphisms of
Proposition~\ref{cut:disk-spatial}, the normalized flag-area measures converge
weakly to the quantum area measure.
\end{proposition}
\begin{proof}\leavevmode

Take \(P=B_n\) and normalize triangle counts by \(2n\), at first even without requiring the total.
At perimeters tending to \(s>0\) in these units the volume law tends by \eqref{eq:original-9} to \(d_s\).
Jointly with the coupled path data in an unconditioned pair, matched cell volumes converge to the areas of the corresponding field disks.
Indeed conditional on parent paths/perimeters the jump fillings are fresh.
Through any fixed finite depth first keep the paths down to the previous depth, and also volumes of cells born at that depth; \eqref{eq:original-9} and the product conditional laws give the stated comparison for the latter jointly by finite tests and exhaustion.
Limiting ancestor masses must dominate the sums of descendant masses by actual containment.
Those sums in the corresponding full continuum generations have exactly the ancestor-area law by area additivity in the physical model; the lattice ancestor marginal already has that same limit by freshness at its own birth and \eqref{eq:original-9}.
Therefore domination is equality (condition on/exhaust ancestor birth tests when selecting ancestors).
Increasing finite depths proves the full joint comparison.
In detail take a subsequential limit with all matched lengths/paths and cell masses jointly (each fixed mass label tight by the perimeter and freshness estimates).
Conditional on paths through any fixed preceding generation the born masses at the next depth then have the continuum product marginals by testing on earlier observation cylinders only; all labels can first be selected on compact ancestry/birth events which do not test their own subsequent interiors.
Thus for any fixed ancestor the sum of these next-generation masses below it already has marginal equal to the true ancestor-mass law by area additivity there, conditional at least on its earlier birth data.
The limiting ancestor variable has this same marginal and dominates the sum, forcing equality a.s., simultaneously through countable levels.
Given any finite joint test keeping also later paths, pass to depth strictly beyond those viewed paths, where the mass marginals conditional on earlier paths are known as above and all earlier masses are the a.s. sums of the indicated descendants.
This is why no independence between mass and the future exploration of its own cell was required.
This works for compatible routings under these area-additive and marked Markov hypotheses.
Consequently normalized flag-area measures converge weakly under the homeomorphisms: distribute according to matched cells of small limiting mesh capturing arbitrarily much of total mass.
Each Tutte triangle accounts for two flags and the initial ring mass is negligible.

\end{proof}

\subsection{Exact sphere volume and removal of the counting bias}\label{cut:conditioning-section}

\begin{proposition}[Transfer to exact total volume]\label{cut:fixed-volume}
Fix a compact positive window for the rescaled perimeter of a centrally
marked loop.  The exact fixed-size counting intensity converges to the
unit-area critical sphere Palm law of Proposition~\ref{cut:palm}, jointly
with every finite compatible surviving path test and every finite collection
of cell-volume labels.  The spatial, area, and full loop-tour comparisons of
Propositions~\ref{cut:disk-spatial} and~\ref{cut:joint-volumes} hold in this
conditioned coupling.
\end{proposition}
\begin{proof}\leavevmode

We detail this step since absolute continuity of a macroscopic finite sphere relative to a pair of Boltzmann disks should not be asserted after exact volume conditioning. Use counting intensity with a central marked loop and restrict \(k/P\in[\alpha,\beta]\Subset(0,\infty)\). Initially on the unconditioned pair fillings put base mass per \((k,l)\)
\[
 b_n(k,l)=\frac{4}{k+l}
 \frac{\binom{k+l}k 2^{-k-l} f_k f_l}{X^{2n}Z_n(4)} .
\]
Use a distinguished link (uniform in the ring in the loop-marked formulation). To obtain the actual restricted marked intensity multiply by \(n{\bf1}_{\{N_k+N_l+k+l=2n\}}\), as in \eqref{eq:original-3}; one may restore an ordinary FK root uniformly with the corresponding occurrence symmetry. Before multiplication these are finite measures with the continuum abstract limit
\[
 \frac{1}{2\pi^2}s^{-5}{\bf1}_{[\alpha,\beta]}(s)\,ds\ 
        \mathrm{Pair}^{\rm prob}(s)
\]
including common-typical boundary alignment, by the same Riemann-sum estimates in \eqref{eq:original-3}, \eqref{eq:original-12} and the volume comparisons.
Here the pair uses the two critical probability disks at perimeter \(s\), with recursive independent signed exploration in the nested coupling, and starts without imposing a total volume.
We need only abstract path/volume labels in this convergence first.
The views after volume conditioning will be coupled to actual \emph{embedded} placements from the continuum Palm law by disintegration, not by assuming conformal-welding uniqueness.

\proofstage{Integrating out one unexamined filling}

Fix a finite compatible path test and finitely many cell-volume probes.
Choose an additional jump filling below the interiors inspected by these
tests.  The choice uses only its birth path, perimeter, and attachment
labels.  In particular it does not inspect the filling itself.  Localize
the choice to a success event on which its degree $p_n$ satisfies
\[
 u_n:=p_n/P\in[u_-,u_+]\Subset(0,\infty).
\]
Compact survival, strict-margin jump matching, and jump density provide a
countable family of such success events exhausting the limiting pair law.
One obtains them by continuing below the finitely many viewed depths until
an additional matched positive-size jump is born.  The selection can use
lengths and times relative to the central perimeter $s$, so the same
exhaustion applies after the scale-integrated area-one conditioning.

Let $\mathcal F_{\rm out}^{(n)}$ contain the selected birth data and all
discrete data outside this filling.  Volume probes of ancestors containing
it are not included as extra conditioning: write each such volume as its
outside contribution plus the reservoir volume, and make the substitution
below.  By the exact peeling factorization, conditional on
$\mathcal F_{\rm out}^{(n)}$ the reservoir is a fresh Boltzmann disk $T_{p_n}$.
Write $N_{\rm res}$ for its triangle count and $N_{\rm out}$ for all other
triangles, including the central ring.  Put
\[
 A_{\rm out}^{(n)}=N_{\rm out}/(2n).
\]
The conditional expectation of the exact-volume multiplier is therefore
\begin{align*}
 &\E\left[n\ind_{\{N_{\rm out}+N_{\rm res}=2n\}}
                           \mid\mathcal F_{\rm out}^{(n)}\right]\\
 &\hspace{2cm}=n\Prob\{N_{p_n}=2n(1-A_{\rm out}^{(n)})\}
   =d_{u_n}(1-A_{\rm out}^{(n)})+o(1).
\end{align*}
The last equality is uniform on the success event by
Lemma~\ref{disk:volume-llt}.  Its parity condition is automatic: the outside
surface with its degree-$p_n$ hole has precisely the complementary triangle
parity.  The density is zero for nonpositive residual area and is uniformly
bounded when $u_n\in[u_-,u_+]$.

All outside quantities converge jointly before conditioning, by
Proposition~\ref{cut:joint-volumes}; for example $A_{\rm out}^{(n)}$ is the
total mass minus the reservoir mass.  In a tested ancestor volume, replace
the reservoir contribution by $1-A_{\rm out}^{(n)}$.  Unrelated volume probes
are unchanged.  The exact conditional identity above now proves convergence
of every bounded continuous test of the retained path data and substituted
volumes on this success event.  It proves the same assertion with additional
bounded outside readouts whenever their unconditioned joint convergence is
available.  It does not require a limit theorem for the interior shape of a
disk conditioned on one exact triangle count.

\proofstage{Identifying the area-one version}

The resulting density selects the area-one Palm law of
Proposition~\ref{cut:palm}, with its specified disintegration.  To verify
this without choosing a conditional law at a null event, start from the
unit-perimeter pair and integrate the central perimeter $s$.  Use bars for
areas in the unit-perimeter pair.  The reservoir has relative perimeter
$u/s$ and, conditional on its outside data, area density $d_{u/s}$.
Scaling the pair to total area one and changing the reservoir-area variable
back to $s$ gives the factor
\[
 s^{-2}d_{u/s}(s^{-2}-\bar A_{\rm out})
       =d_u(1-s^2\bar A_{\rm out})
\]
against the central perimeter measure $s^{-5}\,ds$.  This is exactly the
limiting multiplier above.  Volume entries containing the reservoir undergo
the same substitution.  The calculation remains valid for Borel success
tests chosen from the outside data.

Under this scale-integrated area-one law the unit-perimeter pair variables
are absolutely continuously weighted.  They therefore retain the
almost-sure cutting, contact, and mesh properties of the unconditioned pair.
Finally Proposition~\ref{disk:palm} gives convergence of the whole restricted
counting mass.  Since the reservoir success events exhaust the candidate
area-one law, none of that mass can escape their union.  Refining finite
continuity partitions and then increasing the family of tests proves the
joint assertion for all countably encoded cutting data.  A different,
deeper reservoir may be used for each finite test; no filling is required to
remain independent of the entire limiting tree.

The pinching and tracking estimates used in the spatial comparison also
survive exact volume conditioning.  For a test in one initial disk, keep
the opposite initial disk unexamined and use it as the reservoir.  Both
initial perimeters lie in compact positive ranges, apart from the negligible
ring tails already controlled in Proposition~\ref{disk:palm}.  The uniformly
bounded local-limit multiplier then preserves every vanishing error
probability needed on that test.  The good coupling subsequences required
for Proposition~\ref{cut:disk-spatial} are consequently available under the
exact fixed-size law.

\proofstage{Actual continuum placements and measurable comparisons}

In particular in this restricted counting intensity we have the topological area, nested-loop and length-label comparisons to actual unit sphere cuts.
Conditional on the full matched abstract path and volume data in a subsequential limit one can sample remaining continuum geometry independently of any further lattice readouts by the regular conditional law of the actual field sphere/cuts given those data.
The vertex/incidence consistency and \eqref{eq:original-13} did not require determinism of that conformal sphere from the cutting hierarchy.
All variables used in the coupling may be encoded by countable measurable charts, paths, labels, and probability fields on the standard spherical/disk charts.
Approximate maps and paths, if used probabilistically, can either be chosen measurably by finite triangulation searches up to arbitrarily small additional errors, or just used through the resulting closed sequential comparison assertions.

\end{proof}

\begin{corollary}[The unbiased topological comparison]\label{cut:unbiased}
The ordinary finite FK sphere admits a joint coupling to the prescribed
unit-area critical sphere in which orientation-preserving topological
placements give weak area convergence, full nested-loop matching with
multiplicities, intrinsic loop-length label convergence, and vanishing flag
mesh.  Three independent flag-area samples converge jointly to three
independent quantum-area samples.  The conformal identification of these
placements with the prescribed uniformizations and the comparison of graph
distances are established separately in the following sections.
\end{corollary}
\begin{proof}\leavevmode

To remove Palm bias additionally require the \(m_n\) side fractions of the central loop both to lie above \(\delta>0\), at continuity cutoffs.
In the marked comparison the \emph{entire number} of qualifying loops then converges.
Indeed weak convergence of area to a diffuse measure under the comparison homeomorphisms prevents an arbitrarily small loop from carrying both side fractions bounded below.
All positive-diameter loops are matched one-to-one, with their two degree labels and their volume fractions converging since true loop boundaries have zero area.
Window boundaries can be avoided for the entire countable ensemble, by (the continuum side of) the Palm identity and absolute continuity when re-marking by counting.
The limiting count is finite and at least one on the further restricted law.
The critical lengths for other loops that enter this comparison are the actual intrinsic cut lengths established in the continuum counting Palm identity (hence a.s. for all loops also under central counting restrictions).
For a discovered jump loop its interior chart has precisely the actual pulled-back field there, so by length-as-field-function it has that length with the same units; the two discrete circumference labels agree asymptotically by the ring comparison.
Dividing by the count on both sides now yields the ordinary rooted-map probability restricted to having at least one such loop, coupled to the ordinary unit-area prescribed probability with the same nonemptiness restriction, with convergence of the masses too.
This retains a central choice uniformly among qualifying loops when desired.
Letting the windows increase and \(\delta\) decrease exhausts the latter ordinary probability by positivity/finite length and nonzero side areas of true loops; hence there is no leftover mass in the ordinary map passage either.
One can for example partition by the first nonempty such window.
Signs, if used, correspond to the fair bSLE orientation assignment, while blue and red represent opposite parities from the specified side of the central loop.
Three independent \(m_n\) samples can jointly approach fresh independent area samples by the weak area comparison.
One may postcompose the topological placements accordingly to use the specified orientation-preserving three-mark convention in the continuum.
This proves at this stage the area/CLE/mesh assertions in \emph{coupled topological} positions; replacing these by the actual \(\phi_n\), and comparing intrinsic distances, still requires separate arguments.

\end{proof}
